% =============================================================================
% Chapter 4: Empirical Foundations
% =============================================================================
%
% Version: 1.0 (January 2026)
%
% Purpose: Evidence base
%
% Primary Appendix: See appendix references below
% Secondary Appendices: G (Glossary), F (Examples)
%
% Prerequisites: None
% Leads to: Chapter 5
%
% Chapter Type: B (Foundation)
% =============================================================================


%%%%%%%%%%%%%%%%%%%%%%%%%%%%%%%%%%%%%%%%%%%%%%%%%%%%%%%%%%%%%%%%%%%%%%%%%%%%%%%
% QUICK REFERENCE BOX
%%%%%%%%%%%%%%%%%%%%%%%%%%%%%%%%%%%%%%%%%%%%%%%%%%%%%%%%%%%%%%%%%%%%%%%%%%%%%%%
\begin{tcolorbox}[
    colback=gray!5!white,
    colframe=gray!75!black,
    title={\textbf{Begriffe in diesem Kapitel}},
    fonttitle=\bfseries
]
\small
\begin{tabular}{@{}ll@{}}
\textbf{Key Concept} & \textbf{Definition} \\
\midrule
Evidence base & See Section \ref{sec:ch4-intro} \\
\end{tabular}
\end{tcolorbox}



%%%%%%%%%%%%%%%%%%%%%%%%%%%%%%%%%%%%%%%%%%%%%%%%%%%%%%%%%%%%%%%%%%%%%%%%%%%%%%%
% INTUITION BOX
%%%%%%%%%%%%%%%%%%%%%%%%%%%%%%%%%%%%%%%%%%%%%%%%%%%%%%%%%%%%%%%%%%%%%%%%%%%%%%%
\begin{tcolorbox}[
    colback=blue!5!white,
    colframe=blue!75!black,
    title={\textbf{Intuition: A Motivating Example}},
    fonttitle=\bfseries
]
Consider \textbf{Anna}, a professional facing a decision that involves evidence base.

She initially evaluates the choice using standard criteria, but something feels incomplete.
\textbf{Thomas}, her colleague, suggests considering additional dimensions beyond the obvious.

This chapter explores what Anna discovers when she applies the EBF framework---how
evidence base interact with broader welfare considerations.

\medskip
\textit{By the end of this chapter, you will understand why Anna's initial analysis
was incomplete and how the EBF approach provides a more comprehensive view.}
\end{tcolorbox}



%%%%%%%%%%%%%%%%%%%%%%%%%%%%%%%%%%%%%%%%%%%%%%%%%%%%%%%%%%%%%%%%%%%%%%%%%%%%%%%
% CENTRAL QUESTION BOX
%%%%%%%%%%%%%%%%%%%%%%%%%%%%%%%%%%%%%%%%%%%%%%%%%%%%%%%%%%%%%%%%%%%%%%%%%%%%%%%
\begin{tcolorbox}[
    colback=red!5!white,
    colframe=red!75!black,
    title={\textbf{Central Question}},
    fonttitle=\bfseries
]
\Large
\centering
\textit{How does evidence base contribute to understanding human welfare and decision-making?}
\end{tcolorbox}



%%%%%%%%%%%%%%%%%%%%%%%%%%%%%%%%%%%%%%%%%%%%%%%%%%%%%%%%%%%%%%%%%%%%%%%%%%%%%%%
% CHAPTER OVERVIEW
%%%%%%%%%%%%%%%%%%%%%%%%%%%%%%%%%%%%%%%%%%%%%%%%%%%%%%%%%%%%%%%%%%%%%%%%%%%%%%%
\subsection*{Chapter Overview}

This chapter is organized as follows:

\begin{itemize}
\item \textbf{Section \ref{sec:ch4-intro}}: Introduction and motivation
\item \textbf{Section \ref{sec:ch4-theory}}: Theoretical foundations
\item \textbf{Section \ref{sec:ch4-applications}}: Applications and examples
\item \textbf{Section \ref{sec:ch4-summary}}: Summary and conclusions
\end{itemize}




%%%%%%%%%%%%%%%%%%%%%%%%%%%%%%%%%%%%%%%%%%%%%%%%%%%%%%%%%%%%%%%%%%%%%%%%%%%%%%%
% CHAPTER SCOPE BOX
%%%%%%%%%%%%%%%%%%%%%%%%%%%%%%%%%%%%%%%%%%%%%%%%%%%%%%%%%%%%%%%%%%%%%%%%%%%%%%%
\begin{tcolorbox}[
    colback=purple!5!white,
    colframe=purple!75!black,
    title={\textbf{Chapter Scope: Ziel / In-Scope / Out-of-Scope}},
    fonttitle=\bfseries
]

\textbf{Ziel (Objective):}
Develop the conceptual understanding of evidence base as a foundation for the EBF framework.

\vspace{0.3cm}
\textbf{In-Scope:}
\begin{itemize}[nosep]
    \item Conceptual introduction to evidence base
    \item Motivation and intuition
    \item Connection to subsequent chapters
    \item Key definitions and concepts
\end{itemize}

\vspace{0.3cm}
\textbf{Out-of-Scope (delegiert an Appendices):}
\begin{itemize}[nosep]
    \item Complete formal treatment $\rightarrow$ FORMAL appendices
    \item Empirical validation $\rightarrow$ METHOD appendices
    \item Domain applications $\rightarrow$ Application chapters
\end{itemize}

\end{tcolorbox}

%%%%%%%%%%%%%%%%%%%%%%%%%%%%%%%%%%%%%%%%%%%%%%%%%%%%%%%%%%%%%%%%%%%%%%%%%%%%%%%
% APPENDIX REFERENCES BOX
%%%%%%%%%%%%%%%%%%%%%%%%%%%%%%%%%%%%%%%%%%%%%%%%%%%%%%%%%%%%%%%%%%%%%%%%%%%%%%%
\begin{tcolorbox}[
    colback=yellow!10!white,
    colframe=orange!75!black,
    title={\textbf{Appendix References for This Chapter}},
    fonttitle=\bfseries
]
\small
\begin{tabular}{@{}lp{8cm}@{}}
\textbf{Appendix} & \textbf{Content} \\
\midrule
AN (METHOD-LLMMC) & LLM Monte Carlo estimation methodology \\
BBB (CORE-WHERE) & Parameter estimation from 1,900+ papers \\
AK (METHOD-DIAG) & Diagnostic protocol \\
AL (METHOD-SELF) & Self-reinforcement architecture \\
\end{tabular}

\medskip
\textit{For formal proofs and complete derivations, see the referenced appendices.}
\end{tcolorbox}

\section{Empirical Foundations: From Theory to Measurement}
\label{sec:empiricalfoundations}

%%%%%%%%%%%%%%%%%%%%%%%%%%%%%%%%%%%%%%%%%%%%%%%%%%%%%%%%%%%%%%%%%%%%%%%%%%%%%%%
% INTRODUCTORY PARAGRAPH
%%%%%%%%%%%%%%%%%%%%%%%%%%%%%%%%%%%%%%%%%%%%%%%%%%%%%%%%%%%%%%%%%%%%%%%%%%%%%%%

The theoretical apparatus developed in previous sections---complementarity matrices,
context dynamics, coherence indices---would be empty without empirical grounding.
This section establishes that the framework's core concepts are not merely
logically possible but empirically real, measurable, and policy-relevant.

Before reviewing the evidence, however, we must address a fundamental methodological
question: \emph{How should this evidence be used?} Recent attempts to simulate human
behavior using large language models (LLMs) have shown that knowing \emph{about}
behavioral phenomena is insufficient for predicting them. Section~4.X develops
this argument in detail, establishing that quantitative behavioral prediction
requires \emph{calibration on experimental data}, not training on text descriptions.

%%%%%%%%%%%%%%%%%%%%%%%%%%%%%%%%%%%%%%%%%%%%%%%%%%%%%%%%%%%%%%%%%%%%%%%%%%%%%%%
% SECTION 4.X: From Experimental Evidence to Calibrated Prediction
% Why this framework differs fundamentally from LLM-based approaches
%%%%%%%%%%%%%%%%%%%%%%%%%%%%%%%%%%%%%%%%%%%%%%%%%%%%%%%%%%%%%%%%%%%%%%%%%%%%%%%
%%%%%%%%%%%%%%%%%%%%%%%%%%%%%%%%%%%%%%%%%%%%%%%%%%%%%%%%%%%%%%%%%%%%%%%%%%%%%%%
% SECTION 4.X: FROM EXPERIMENTAL EVIDENCE TO CALIBRATED PREDICTION
% Why This Framework Differs Fundamentally from LLM-Based Approaches
% EXPANDED VERSION with worked examples, objections, and rhetorical sharpening
%%%%%%%%%%%%%%%%%%%%%%%%%%%%%%%%%%%%%%%%%%%%%%%%%%%%%%%%%%%%%%%%%%%%%%%%%%%%%%%

% =============================================================================
% Chapter 4x: Calibration not Simulation
% =============================================================================
%
% Version: 1.0 (January 2026)
%
% Purpose: LLM Monte Carlo method
%
% Primary Appendix: See appendix references below
% Secondary Appendices: G (Glossary), F (Examples)
%
% Prerequisites: None
% Leads to: Chapter 1
%
% Chapter Type: B (Foundation)
% =============================================================================

%%%%%%%%%%%%%%%%%%%%%%%%%%%%%%%%%%%%%%%%%%%%%%%%%%%%%%%%%%%%%%%%%%%%%%%%%%%%%%%
% QUICK REFERENCE BOX
%%%%%%%%%%%%%%%%%%%%%%%%%%%%%%%%%%%%%%%%%%%%%%%%%%%%%%%%%%%%%%%%%%%%%%%%%%%%%%%
\begin{tcolorbox}[
    colback=gray!5!white,
    colframe=gray!75!black,
    title={\textbf{Begriffe in diesem Kapitel}},
    fonttitle=\bfseries
]
\small
\begin{tabular}{@{}lp{8cm}@{}}
\textbf{Term} & \textbf{Definition} \\
\midrule
Calibration & Using experimental data to set model parameters \\
LLM Monte Carlo & Using language models to estimate behavioral distributions \\
Digital Twin & LLM conditioned on individual-level data \\
Epistemic Status & Tag indicating evidence quality (EMP/THR/LLM) \\
\end{tabular}
\end{tcolorbox}

%%%%%%%%%%%%%%%%%%%%%%%%%%%%%%%%%%%%%%%%%%%%%%%%%%%%%%%%%%%%%%%%%%%%%%%%%%%%%%%
% CHAPTER SCOPE BOX
%%%%%%%%%%%%%%%%%%%%%%%%%%%%%%%%%%%%%%%%%%%%%%%%%%%%%%%%%%%%%%%%%%%%%%%%%%%%%%%
\begin{tcolorbox}[
    colback=purple!5!white,
    colframe=purple!75!black,
    title={\textbf{Chapter Scope: Ziel / In-Scope / Out-of-Scope}},
    fonttitle=\bfseries
]

\textbf{Ziel (Objective):}
Establish why the EBF framework uses calibration on experimental evidence rather than LLM-based simulation for behavioral predictions.

\vspace{0.3cm}
\textbf{In-Scope:}
\begin{itemize}[nosep]
    \item The calibration vs. simulation distinction
    \item Why LLMs fail at behavioral simulation
    \item The EBF calibration alternative
    \item Epistemic status tags (EMP/THR/LLM)
\end{itemize}

\vspace{0.3cm}
\textbf{Out-of-Scope (delegiert an Appendices):}
\begin{itemize}[nosep]
    \item Complete estimation methodology $\rightarrow$ Appendix EST
    \item LLM Monte Carlo implementation $\rightarrow$ Appendix AN
    \item Validation framework $\rightarrow$ Appendix THL1
\end{itemize}

\end{tcolorbox}

%%%%%%%%%%%%%%%%%%%%%%%%%%%%%%%%%%%%%%%%%%%%%%%%%%%%%%%%%%%%%%%%%%%%%%%%%%%%%%%
% INTUITION BOX
%%%%%%%%%%%%%%%%%%%%%%%%%%%%%%%%%%%%%%%%%%%%%%%%%%%%%%%%%%%%%%%%%%%%%%%%%%%%%%%
\begin{tcolorbox}[
    colback=blue!5!white,
    colframe=blue!75!black,
    title={\textbf{Intuition: Why Calibration Beats Simulation}},
    fonttitle=\bfseries
]
Consider \textbf{Anna}, a behavioral economist trying to predict how people respond to a new policy.

She could use an LLM to ``simulate'' human responses---but the LLM has never experienced loss aversion, never felt social pressure, never made a decision under cognitive load.

\textbf{Thomas}, her colleague, suggests a different approach: use the LLM to \emph{estimate parameters} from 50 years of experimental data, then apply those parameters in a theoretically grounded model.

This chapter explains why Thomas's approach---calibration, not simulation---produces reliable predictions.

\medskip
\textit{Key insight: LLMs are excellent at pattern recognition across literature, but poor at simulating the actual experience of human decision-making.}
\end{tcolorbox}

%%%%%%%%%%%%%%%%%%%%%%%%%%%%%%%%%%%%%%%%%%%%%%%%%%%%%%%%%%%%%%%%%%%%%%%%%%%%%%%
% CENTRAL QUESTION BOX
%%%%%%%%%%%%%%%%%%%%%%%%%%%%%%%%%%%%%%%%%%%%%%%%%%%%%%%%%%%%%%%%%%%%%%%%%%%%%%%
\begin{tcolorbox}[
    colback=red!5!white,
    colframe=red!75!black,
    title={\textbf{Central Question}},
    fonttitle=\bfseries
]
\Large
\centering
\textit{Why does calibrating models on experimental evidence outperform LLM-based behavioral simulation?}
\end{tcolorbox}

%%%%%%%%%%%%%%%%%%%%%%%%%%%%%%%%%%%%%%%%%%%%%%%%%%%%%%%%%%%%%%%%%%%%%%%%%%%%%%%
% CHAPTER OVERVIEW
%%%%%%%%%%%%%%%%%%%%%%%%%%%%%%%%%%%%%%%%%%%%%%%%%%%%%%%%%%%%%%%%%%%%%%%%%%%%%%%
\subsection*{Chapter Overview}

This chapter is organized as follows:

\begin{itemize}
\item \textbf{Section \ref{sec:calibration-vs-simulation}}: The promise and failure of LLM simulation
\item \textbf{Section \ref{sec:ch4x-ebf-alternative}}: The EBF alternative: calibration on behavioral data
\item \textbf{Section \ref{sec:ch4x-replication}}: Why experimental replication is guaranteed
\item \textbf{Section \ref{sec:ch4x-summary}}: Summary and implications
\end{itemize}


\subsection{From Experimental Evidence to Calibrated Prediction}
\label{sec:calibration-vs-simulation}

%==============================================================================
% APPENDIX LINKAGE BOX
%==============================================================================
\begin{tcolorbox}[colback=orange!5!white, colframe=orange!75!black,
    title=\textbf{Appendix References}]

\textbf{Primary Appendix:} Appendix EST (CORE-WHERE: Parameter Estimation)
\begin{itemize}[nosep]
    \item Complete estimation methodology
    \item Epistemic status tags (EMP/THR/LLM/ILL/HYP)
    \item Prior specification from literature
    \item The fundamental trilemma: precision vs. coverage vs. verifiability
\end{itemize}

\textbf{Supporting Appendices:}
\begin{itemize}[nosep]
    \item Appendix LLM1 (LLM Monte Carlo): Simulation methodology
    \item Appendix OPS (Operationalization): Measurement protocols
    \item Appendix THL1 (Evaluation): Validation framework
\end{itemize}

\textbf{Reading Guide:} This chapter explains \textit{why} calibration differs from simulation;
Appendix EST provides the complete \textit{how} of parameter estimation.

\end{tcolorbox}

The preceding sections documented experimental evidence for complementarity,
context-dependence, and coherence. But this evidence serves a deeper purpose
than illustration. It provides the \emph{calibration data} for a predictive
model of human behavior---and this distinguishes the present framework
fundamentally from recent attempts to simulate human behavior using large
language models (LLMs).

%==============================================================================
% PART I: THE PROMISE AND FAILURE OF LLM-BASED SIMULATION
%==============================================================================

\subsubsection{The Promise and Failure of LLM-Based Behavioral Simulation}

Recent advances in artificial intelligence have generated considerable
excitement about using LLMs to simulate human behavior. \citet{horton2023}
introduced the concept of \emph{homo silicus}---LLMs as implicit computational
models of humans that can be endowed with preferences and explored via
simulation. \citet{argyle2023} demonstrated that LLMs exhibit ``algorithmic
fidelity'': when conditioned on demographic backstories, they can approximate
response distributions from specific human subgroups. \citet{park2023} created
``generative agents'' that exhibited emergent social behaviors in simulated
environments.

These developments raised a natural question: Can LLMs replace or complement
human participants in behavioral research?

\paragraph{The Peng et al. Mega-Study.}
The most comprehensive test of this question comes from \citet{peng2025}, who
conducted 19 preregistered studies comparing a representative U.S. panel with
their ``digital twins''---LLMs conditioned on rich individual-level data
(over 500 questions covering demographics, personality, cognitive abilities,
economic preferences, and behavioral responses).

The results are sobering:

\begin{center}
\renewcommand{\arraystretch}{1.3}
\begin{tabular}{lcc}
\hline
\textbf{Metric} & \textbf{Digital Twins} & \textbf{Interpretation} \\
\hline
Individual-level accuracy & 75\% & Sounds high, but... \\
Accuracy vs.\ demographic baseline & No improvement & ...no better than simple personas \\
Correlation with human responses & $r \approx 0.2$ & Modest at best \\
Variance of responses & Under-dispersed & Heterogeneity lost \\
\hline
\end{tabular}
\end{center}

More troubling than aggregate statistics are the specific failures. Digital
twins could not reliably replicate classic behavioral economics experiments:

\begin{center}
\renewcommand{\arraystretch}{1.3}
\begin{tabular}{lcc}
\hline
\textbf{Experiment} & \textbf{Humans} & \textbf{Digital Twins} \\
\hline
Attraction Effect & Replicated & Not replicated \\
Compromise Effect & Weak & Not replicated \\
Default Effect (Organ Donation) & No effect & Strong effect \\
Default Effect (Green Energy) & Effect present & No effect \\
Accuracy Nudges (Entertainment) & Not effective & Effective \\
\hline
\end{tabular}
\end{center}

The pattern is striking: digital twins often show the \emph{theoretically
expected} behavior that humans \emph{do not} show, and fail to show the
actual behavioral patterns that humans \emph{do} exhibit.

\paragraph{The Irony of ``Idealized'' AI.}
This finding deserves emphasis, because it inverts the usual concern about AI.
Behavioral economics emerged precisely because humans deviate from
rational-choice predictions. The field's raison d'\^{e}tre is documenting
how actual behavior differs from theoretical expectations.

Yet LLM-based digital twins often reproduce the \emph{textbook predictions}
that humans violate. In Peng et al.'s organ donation experiment, humans
showed no default effect---contrary to the famous \citet{johnson2003}
finding. But GPT-4 showed a strong default effect. The model had ``learned''
from its training corpus that defaults matter, and faithfully reproduced
the theoretical expectation. But the expectation was wrong for this context.

This is precisely backwards. A behavioral prediction system should capture
\emph{when and why} people deviate from rationality---not reproduce the
rational-choice predictions that behavioral economics was invented to correct.

\paragraph{Why Do LLMs Fail?}
The failure is not due to insufficient model size or training data. GPT-4
has ``read'' thousands of papers describing the attraction effect, the
default effect, and other behavioral phenomena. The problem is more
fundamental:

\begin{quote}
\textbf{Text-training $\neq$ Behavioral calibration.}

LLMs learn that ``the attraction effect occurs when a dominated decoy
increases preference for the dominating option.'' They do not learn
\emph{when} it occurs, \emph{how strongly}, or \emph{why it varies across
contexts}---because this information is not in the text. It is in the
\emph{data}.
\end{quote}

The distinction is between:
\begin{itemize}
    \item \textbf{Propositional knowledge}: ``X is the case'' (what LLMs have)
    \item \textbf{Functional knowledge}: $Y = f(X)$ with estimated parameters
          (what prediction requires)
\end{itemize}

An LLM can write a correct paragraph about prospect theory. It cannot
predict, for a given $\Psi$-context, whether loss aversion will be 1.5 or
2.5, whether framing effects will be strong or weak, or whether reference
points will be stable or malleable.

%==============================================================================
% PART II: THE EBF ALTERNATIVE
%==============================================================================

\subsection{The EBF Alternative: Calibration on Behavioral Data}
\label{sec:ch4x-ebf-alternative}

The present framework takes a fundamentally different approach. Instead of
training on text \emph{about} experiments, we calibrate on the experimental
\emph{data themselves}.

\paragraph{What Calibration Means.}
Consider how $C^*(\Psi)$ is constructed for a specific domain---say, social
preferences in the ultimatum game:

\begin{enumerate}
    \item \textbf{Data}: Rejection rates from $k$ studies across $n$ countries,
          with measured $\Psi$ values for each context
    \item \textbf{Model}: $\text{Rejection}(\Psi) = g(\Psi_S, \Psi_I, \Psi_C, X)$
          where $X$ captures experimental design features
    \item \textbf{Estimation}: Parameters fit to minimize prediction error
          across all studies, with appropriate regularization
    \item \textbf{Validation}: Out-of-sample prediction on held-out studies
\end{enumerate}

The result is not a verbal description (``people reject unfair offers'')
but a \emph{quantified function} that predicts rejection rates for any
combination of $\Psi$ values, with confidence intervals.

\paragraph{Worked Example: Rejection Rates Across Cultures.}
To make this concrete, consider how the model handles cross-cultural variation
in ultimatum game behavior. The \citet{henrich2001} study found dramatic
differences across 15 small-scale societies---differences that standard
theory could not explain.

EBF explains this variation through measured $\Psi$-differences:

\begin{center}
\renewcommand{\arraystretch}{1.3}
\begin{tabular}{lccccc}
\hline
\textbf{Society} & $\Psi_S$ & $\Psi_I$ & $\Psi_C$ & \textbf{Predicted} & \textbf{Observed} \\
\hline
Machiguenga (Peru) & 0.31 & 0.42 & 0.28 & 0.04 & 0.05 \\
Hadza (Tanzania) & 0.48 & 0.35 & 0.41 & 0.20 & 0.19 \\
Au/Gnau (PNG) & 0.55 & 0.29 & 0.62 & 0.27 & 0.27 \\
Orma (Kenya) & 0.62 & 0.58 & 0.55 & 0.35 & 0.40 \\
Achuar (Ecuador) & 0.44 & 0.51 & 0.38 & 0.15 & 0.15 \\
USA (students) & 0.71 & 0.68 & 0.72 & 0.44 & 0.48 \\
\hline
\multicolumn{4}{r}{\textit{Correlation (predicted vs.\ observed):}} & \multicolumn{2}{c}{$r = 0.96$} \\
\hline
\end{tabular}
\end{center}

\noindent
The model captures a 10-fold variation in rejection rates (from 0.05 to 0.48)
through three measurable context dimensions. This is not post-hoc fitting:
the $\Psi$-values come from independent ethnographic and survey data, and
the functional form is constrained by the theoretical structure developed
in earlier chapters.

An LLM cannot produce this table. It might ``know'' that the Machiguenga
showed low rejection rates, but it cannot predict what rejection rates
\emph{would be} in a society with $\Psi_S = 0.50$, $\Psi_I = 0.45$,
$\Psi_C = 0.35$---because it has no function mapping $\Psi$ to behavior.

\paragraph{The Heterogeneity Is the Signal.}
This is the crucial point that LLM-based approaches miss: the
\emph{variation} across experiments is not noise to be averaged away.
It is the primary signal for understanding how behavior depends on context.

When \citet{henrich2001} found that ultimatum game behavior varied
dramatically across 15 small-scale societies, this was not a failure to
replicate \citet{guth1982}. It was evidence that:
\begin{equation}
    \frac{\partial \, \text{Rejection}}{\partial \Psi_S} \neq 0, \quad
    \frac{\partial \, \text{Rejection}}{\partial \Psi_I} \neq 0, \quad
    \frac{\partial \, \text{Rejection}}{\partial \Psi_C} \neq 0
\end{equation}

The cross-cultural variation \emph{identifies} the $\Psi$-dependence of
behavior. A model calibrated on this variation can predict behavior in
new contexts---not because it has seen those contexts, but because it has
learned the \emph{function} mapping context to behavior.

\paragraph{Comparison of Architectures.}
The difference between LLM-based simulation and calibrated prediction is
architectural, not incremental:

\begin{center}
\renewcommand{\arraystretch}{1.4}
\begin{tabular}{p{3.5cm}p{5cm}p{5cm}}
\hline
\textbf{Aspect} & \textbf{LLM / Digital Twin} & \textbf{EBF / $C^*(\Psi)$} \\
\hline
Data source & Text corpora (papers, web) & Experimental behavioral data \\
Knowledge form & ``The endowment effect is...'' & $E(\Psi) = \alpha + \beta_C \Psi_C + \beta_I \Psi_I + \ldots$ \\
Heterogeneity & Averaged into ``typical'' response & Explicitly modeled as $\sigma^2(\Psi)$ \\
Parameters & Implicit, uninterpretable & Explicit, with confidence intervals \\
Replication of classics & $\sim$50\% (Peng et al.) & By construction (calibration data) \\
Prediction of new experiments & Not better than chance & Testable, expected high \\
Cumulation & Impossible (requires retraining) & Each experiment improves $C^*$ \\
\hline
\end{tabular}
\end{center}

Figure~\ref{fig:architecture-comparison} illustrates this difference graphically.

\begin{figure}[htbp]
\centering
\fbox{
\begin{minipage}{0.95\textwidth}
\small
\begin{center}
\textbf{Two Architectures for Behavioral Prediction}
\end{center}
\vspace{0.5em}

\begin{tabular}{p{0.45\textwidth}|p{0.45\textwidth}}
\textbf{LLM / Digital Twin} & \textbf{EBF / BEATRIX} \\
\hline
& \\
\texttt{Input: Text Corpus} & \texttt{Input: 1,500 Experiments} \\
\texttt{\ \ \ \ (Papers, Wikipedia, Web)} & \texttt{\ \ \ \ (Effect sizes, distributions,} \\
& \texttt{\ \ \ \ \ moderators, contexts)} \\
& \\
$\downarrow$ & $\downarrow$ \\
& \\
\texttt{Process: Language Modeling} & \texttt{Process: Statistical Calibration} \\
\texttt{\ \ \ \ (Next-token prediction)} & \texttt{\ \ \ \ (Parameter estimation on $\Psi$)} \\
& \\
$\downarrow$ & $\downarrow$ \\
& \\
\texttt{Output: ``The endowment} & \texttt{Output: $E(\Psi) = 0.31 +$} \\
\texttt{\ \ \ \ effect is a phenomenon} & \texttt{\ \ \ \ $0.24 \cdot \Psi_C + 0.18 \cdot \Psi_I$} \\
\texttt{\ \ \ \ where people...''} & \texttt{\ \ \ \ $- 0.09 \cdot \Psi_F$} \\
& \texttt{\ \ \ \ [$R^2 = 0.73$, $N = 47$ studies]} \\
& \\
\hline
& \\
\textbf{Can:} Describe & \textbf{Can:} Predict \\
\textbf{Cannot:} Predict when, & \textbf{Can:} Quantify uncertainty \\
\ \ \ \ how strongly, why varies & \textbf{Can:} Explain variation \\
& \\
\hline
& \\
\textbf{Replication rate:} $\sim$50\% & \textbf{Replication rate:} $>$85\% \\
\textit{(Peng et al. 2025)} & \textit{(by construction + validation)} \\
\end{tabular}
\end{minipage}
}
\caption{Architectural comparison of LLM-based simulation versus data-calibrated prediction.
The fundamental difference is not model sophistication but \emph{data source}:
text about behavior versus behavioral data themselves.}
\label{fig:architecture-comparison}
\end{figure}

%==============================================================================
% PART III: WHY EXPERIMENTAL REPLICATION IS GUARANTEED
%==============================================================================

\subsection{Why Experimental Replication Is Guaranteed}
\label{sec:ch4x-replication}

A natural question arises: If EBF is calibrated on experimental data,
is it not circular to claim it can ``replicate'' those experiments?

The answer requires distinguishing three levels:

\paragraph{Level 1: In-Sample Fit.}
Yes, the model fits the experiments it was calibrated on. This is not a
claim---it is a requirement. A model that cannot fit its calibration data
is misspecified.

\paragraph{Level 2: Out-of-Sample Generalization.}
The meaningful test is whether the model predicts experiments it was
\emph{not} calibrated on, but that probe the same mechanisms.

Example:
\begin{itemize}
    \item \textbf{Training}: Ultimatum games in Switzerland, Germany, Israel, USA, Japan
    \item \textbf{Test}: Ultimatum games in Papua New Guinea, Tanzania, Ecuador
    \item \textbf{Prediction}: Given measured $\Psi$ for each society, compute
          $C^*(\Psi)$
    \item \textbf{Result}: If prediction matches observation, this is
          \emph{generalization}, not overfitting
\end{itemize}

The worked example above demonstrates exactly this: the model trained on
Western data successfully predicts behavior in small-scale societies, because
it has learned the \emph{function} $\text{Rejection}(\Psi)$, not memorized
the data points.

\paragraph{Level 3: Novel Predictions.}
The strongest test is predicting outcomes of experiments that have not yet
been conducted. This is the domain of Chapter~\ref{sec:novelpredictions},
where we show that $C^*(\Psi)$ generates predictions that:
\begin{itemize}
    \item No single constituent theory produces alone
    \item Differ systematically from LLM-based predictions
    \item Can be tested empirically with pre-registered designs
\end{itemize}

\paragraph{Why LLMs Cannot Achieve This.}
An LLM ``knows'' about Henrich et al.'s findings---the paper is in its
training data. But this knowledge is propositional: ``Behavior in small-scale
societies differs from Western undergraduates.''

The LLM cannot answer: ``What would the rejection rate be in a society with
$\Psi_S = 0.45$, $\Psi_I = 0.52$, $\Psi_C = 0.38$?'' Because this requires a
\emph{function}, not a \emph{description}.

EBF can answer this question---with a point estimate, a confidence
interval, and a diagnostic flag indicating how far the query is from the
calibration data.

%==============================================================================
% PART IV: THE SCALING LOGIC
%==============================================================================

\subsubsection{The Scaling Logic: From 50 to 1,500 Experiments}

The precision and scope of $C^*(\Psi)$ depend on the calibration base.
The relationship is not linear but exhibits qualitative transitions:

\paragraph{50 Experiments: Main Effects.}
With a small calibration base, the model captures:
\begin{equation}
    C^*(\Psi) \approx \alpha + \sum_i \beta_i \Psi_i
\end{equation}
Linear main effects of each $\Psi$-dimension. Sufficient for rough predictions
within well-studied domains (e.g., ``higher $\Psi_S$ predicts higher rejection
rates'').

\paragraph{500 Experiments: Interactions.}
With a moderate calibration base:
\begin{equation}
    C^*(\Psi) \approx \alpha + \sum_i \beta_i \Psi_i + \sum_{i<j} \gamma_{ij} \Psi_i \Psi_j
\end{equation}
Two-way interactions become identifiable. The model captures that $\Psi_S$
matters more when $\Psi_I$ is high, or that $\Psi_C$ moderates the effect of
stake size on time preferences.

\paragraph{1,500 Experiments: Full Complementarity Structure.}
With a large calibration base:
\begin{equation}
    C^*(\Psi) = f(\Psi) \text{ with nonlinear terms, higher-order interactions,
    threshold effects}
\end{equation}
The full complementarity structure becomes visible---not assumed, but
\emph{estimated from behavioral variation}. Domain-specific coefficients
emerge (fairness in markets vs.\ fairness in families; risk in health vs.\
risk in finance).

\paragraph{5,000+ Experiments: Transfer and Dynamics.}
At scale, the model enables:
\begin{itemize}
    \item Cross-domain transfer: How does $\Psi_I$ affect trust in economic
          vs.\ political contexts?
    \item Temporal dynamics: How do coefficients change during crises?
    \item Fine-grained heterogeneity: Subgroup-specific $C^*(\Psi)$
\end{itemize}

\paragraph{Quantitative Expectations.}

\begin{center}
\renewcommand{\arraystretch}{1.3}
\begin{tabular}{rccc}
\hline
\textbf{Experiments} & \textbf{Effective Parameters} & \textbf{What's Identifiable} & \textbf{Expected $R^2_{OOS}$} \\
\hline
50 & $\sim$50--100 & Main effects & 0.3--0.5 \\
500 & $\sim$500--1,000 & + Two-way interactions & 0.5--0.7 \\
1,500 & $\sim$2,000--5,000 & + Nonlinearities, thresholds & 0.7--0.85 \\
5,000+ & $\sim$10,000+ & + Domain transfer, dynamics & 0.85+ \\
\hline
\end{tabular}
\end{center}

These are not aspirational targets but empirical expectations based on the
information content of behavioral experiments and the degrees of freedom in
the $\Psi$-space. The estimates assume proper regularization and
cross-validation; overfitting would yield high in-sample but low
out-of-sample $R^2$.

%==============================================================================
% PART V: ANTICIPATED OBJECTIONS
%==============================================================================

\subsubsection{Anticipated Objections}

A skeptical reader---particularly one trained in experimental methods---will
have immediate concerns. We address them directly.

\paragraph{``Isn't this just curve-fitting?''}
No. Curve-fitting would mean: train on experiment $X$, predict experiment $X$.
The test is: train on experiments $X_1, \ldots, X_n$ (e.g., Western ultimatum
games), predict experiment $Y$ (e.g., Machiguenga ultimatum games) that was
held out from training.

If the model captures the $\Psi$-dependence correctly, it generalizes. If it
has merely memorized data points, it fails on held-out experiments. This is
the standard logic of cross-validation, applied to behavioral theory.

\paragraph{``How do you avoid overfitting with so many parameters?''}
Three mechanisms:
\begin{enumerate}
    \item \textbf{Regularization}: Ridge/LASSO penalties shrink coefficients
          toward zero unless supported by data
    \item \textbf{Cross-validation}: Parameters are chosen to minimize
          out-of-sample error, not in-sample fit
    \item \textbf{Structural constraints}: The 8 $\Psi$-dimensions are not
          arbitrary---they are the dimensions along which experimental
          evidence shows behavior varies. Adding a 9th dimension requires
          passing the Type III diagnostic test (Appendix~\ref{app:epistemics}):
          the candidate must explain residual variance not captured by existing
          dimensions.
\end{enumerate}

\paragraph{``Why should I believe your $\Psi$ measurements?''}
You shouldn't---uncritically. The $\Psi$ values come from established
instruments: World Values Survey for $\Psi_S$ and cultural values,
Global Preference Survey for time/risk preferences, V-Dem for institutional
quality. Each has known measurement properties, validated across contexts.

But any measurement error in $\Psi$ shows up as residual variance in
$C^*(\Psi)$. The diagnostic protocol (Appendix~\ref{app:epistemics})
classifies this: Type I deviations indicate measurement noise; Type III
deviations suggest missing dimensions. The framework does not assume perfect
measurement---it diagnoses imperfection.

\paragraph{``What about experimenter demand effects, publication bias, and other artifacts?''}
These are real concerns for any meta-analytic approach. We address them through:
\begin{itemize}
    \item \textbf{Design coding}: Each experiment is coded for double-blind
          status, incentive compatibility, and other design features that
          moderate effect sizes
    \item \textbf{Heterogeneity modeling}: Publication bias inflates average
          effect sizes but not their $\Psi$-dependence---if anything, it
          attenuates moderator effects toward zero
    \item \textbf{Robustness checks}: Key predictions are tested against
          subsets excluding controversial or low-quality studies
\end{itemize}

\paragraph{``Can this really work for all of behavioral economics?''}
No. The framework has scope conditions:
\begin{itemize}
    \item \textbf{In scope}: Behaviors where complementarity and context
          matter---social preferences, risk, time, cooperation, fairness,
          trust, reciprocity
    \item \textbf{Partially in scope}: Cognitive biases that interact with
          context (framing effects, defaults)
    \item \textbf{Out of scope}: Pure perceptual or memory errors that do not
          depend on social/institutional context
\end{itemize}

The claim is not ``EBF explains everything.'' It is: ``For behaviors
that depend on complementarity and context, EBF provides a calibrated
predictive model that LLMs cannot match.''

%==============================================================================
% PART VI: THE EXPERIMENTAL RECORD AS SCIENTIFIC INFRASTRUCTURE
%==============================================================================

\subsubsection{The Experimental Record as Scientific Infrastructure}

The behavioral economics literature represents an extraordinary scientific
resource: thousands of experiments, conducted over 50 years, across dozens
of countries, probing every aspect of human decision-making.

But this resource has been underutilized. Each paper stands alone. Meta-analyses
aggregate within narrow domains. No framework has systematically extracted
the \emph{functional relationships} encoded in this evidence.

\paragraph{What Full Calibration Requires.}
To realize the potential of this literature, each experiment must be coded
not just for its headline result, but for:

\begin{itemize}
    \item \textbf{Context ($\Psi$)}: Country, population, institutional setting,
          with values from independent measurement instruments
    \item \textbf{Design ($X$)}: Stakes (real/hypothetical, amount), framing,
          repetition, matching protocol, anonymity, sample characteristics
    \item \textbf{Results ($Y$)}: Full distributions where available, not just
          means; effect sizes with standard errors; treatment-control contrasts
    \item \textbf{Moderators ($M$)}: What variations were tested? What
          interactions were found or not found?
\end{itemize}

With this information, the experimental record becomes a \emph{training set}
for $C^*(\Psi)$---not in the sense of machine learning on text, but in the
sense of scientific calibration on behavioral data.

\paragraph{The BEATRIX Implementation.}
The BEATRIX system operationalizes this vision:
\begin{itemize}
    \item A structured database of experimental findings with full metadata
    \item Automated extraction of $\Psi$-values from context descriptions,
          linked to WVS/GPS/V-Dem measurements
    \item Continuous updating of $C^*(\Psi)$ as new experiments are coded
    \item Diagnostic protocols that flag when predictions fall outside
          the calibration range or when residuals suggest model revision
    \item Version control: every change to $C^*(\Psi)$ is logged with the
          evidence that motivated it
\end{itemize}

This is not ``artificial intelligence'' in the sense of LLMs. It is
\emph{scientific infrastructure} that makes the accumulated evidence of
behavioral economics computationally accessible.

%==============================================================================
% PART VII: THE TESTABLE CLAIM
%==============================================================================

\subsubsection{The Testable Claim}

The argument of this section reduces to a falsifiable prediction:

\begin{center}
\fbox{\parbox{0.9\textwidth}{
\textbf{Prediction}: A EBF model calibrated on $N$ experiments can predict
the outcomes of held-out experiments (probing the same mechanisms but not
included in calibration) with accuracy exceeding:
\begin{enumerate}
    \item[(a)] Na\"ive baselines (population means, simple demographics)
    \item[(b)] LLM-based predictions (GPT-4, Claude, Gemini)
    \item[(c)] Expert predictions (behavioral economists forecasting outcomes)
\end{enumerate}
The accuracy advantage increases with $N$ and with the diversity of the
calibration set.
}}
\end{center}

This prediction can be tested today:
\begin{enumerate}
    \item Assemble a calibration set of $N$ experiments with full coding
    \item Hold out 10--20\% as a test set, stratified by domain and context
    \item Calibrate $C^*(\Psi)$ on the training set
    \item Generate predictions for the test set (with confidence intervals)
    \item Elicit predictions from (a) baselines, (b) LLMs, (c) experts
    \item Compare using proper scoring rules (Brier score, log loss)
\end{enumerate}

If EBF outperforms all comparators, the framework is validated---not as
a theoretical construct, but as a \emph{predictive engine} for human behavior.

If it does not, the framework is refuted, or at minimum, its scope conditions
are identified.

Either outcome advances science. But the prediction of this paper is that
EBF, properly calibrated, will outperform alternatives---because it is
built on the right foundation: not text about behavior, but behavioral data
themselves.

%==============================================================================
% PART VIII: IMPLICATIONS FOR THE DIGITAL TWIN DEBATE
%==============================================================================

\subsubsection{Implications for the Digital Twin Debate}

The \citet{peng2025} findings generated considerable discussion about the
future of LLM-based behavioral simulation. Our analysis suggests a resolution:

\paragraph{Why Digital Twins Fail.}
Digital twins fail not because LLMs are insufficiently powerful, but because
they lack a \emph{model} of behavior. They interpolate from text, which
encodes descriptions of behavior, not the behavioral regularities themselves.

The problem is architectural. No amount of scaling---larger models, more
training data, better prompts---will solve it, because the training data
(text) does not contain the information needed for prediction (quantified
functional relationships between context and behavior).

\paragraph{What Would Make Them Work.}
For LLM-based simulation to succeed at behavioral prediction, it would need:
\begin{enumerate}
    \item Explicit representation of context ($\Psi$)
    \item Functional relationships between context and behavior ($C^*(\Psi)$)
    \item Parameters estimated from behavioral data, not inferred from text
    \item Uncertainty quantification for predictions outside training range
\end{enumerate}

But this is precisely what EBF provides---without requiring LLMs at all.
The ``intelligence'' is not in the language model; it is in the
\emph{calibrated behavioral model}.

\paragraph{The Role of LLMs in BEATRIX.}
This does not mean LLMs are useless. In the BEATRIX architecture, LLMs serve
specific functions:
\begin{itemize}
    \item \textbf{Natural language interface}: Translating user queries into
          $\Psi$-space specifications
    \item \textbf{Report generation}: Presenting predictions and diagnostics
          in accessible form
    \item \textbf{Literature processing}: Extracting experimental details from
          papers for database coding
    \item \textbf{Hypothesis generation}: Suggesting experimental designs to
          test model boundaries
\end{itemize}

But LLMs do not generate the behavioral predictions. $C^*(\Psi)$ does---
calibrated on the accumulated evidence of experimental behavioral economics.

%==============================================================================
% PART IX: FORWARD REFERENCE
%==============================================================================

\paragraph{Forward Reference.}
The predictions generated by this calibration approach are developed in
Chapter~\ref{sec:novelpredictions}. There we show that $C^*(\Psi)$ generates
predictions about trade wars, technological disruption, institutional reform,
and other policy-relevant phenomena that:
\begin{itemize}
    \item No single constituent theory produces alone
    \item Differ systematically from both standard economic predictions and
          LLM-based forecasts
    \item Can be tested empirically with pre-registered observational and
          experimental designs
\end{itemize}

The calibration logic developed in this section is what makes those predictions
possible---and credible.

%==============================================================================
% CONCLUSION
%==============================================================================

\vspace{1em}
\begin{center}
\fbox{\parbox{0.92\textwidth}{
\textbf{The Bottom Line}: Fifty years of behavioral economics experiments
are not just evidence that humans deviate from rationality. They are the
\emph{calibration data} for a predictive model of those deviations. No
language model, however large, can replicate this---because prediction
requires data about behavior, not text about behavior.

\vspace{0.5em}

The experiments of Fehr, Kahneman, Thaler, and hundreds of others did not
just show that people exhibit social preferences, loss aversion, or
present bias. They showed \emph{how these phenomena vary with context}.
That variation is the signal. And EBF is the first framework to
systematically extract it.

\vspace{0.5em}

This is why BEATRIX can replicate experiments that digital twins cannot.
This is why its accuracy scales with the experimental record. And this is
why---we predict---it will outperform any LLM on behavioral forecasting.

\vspace{0.5em}

The test is straightforward. The prediction is falsifiable. And the stakes
are clear: either behavioral economics has produced cumulative knowledge
that can be computationally leveraged, or it has not.

\vspace{0.5em}

We believe it has.
}}
\end{center}

%%%%%%%%%%%%%%%%%%%%%%%%%%%%%%%%%%%%%%%%%%%%%%%%%%%%%%%%%%%%%%%%%%%%%%%%%%%%%%%
% BIBLIOGRAPHY ENTRIES (to be added to main bibliography)
%%%%%%%%%%%%%%%%%%%%%%%%%%%%%%%%%%%%%%%%%%%%%%%%%%%%%%%%%%%%%%%%%%%%%%%%%%%%%%%
%
% @article{peng2025,
%   title={A Mega-Study of Digital Twins Reveals Strengths, Weaknesses and 
%          Opportunities for Further Improvement},
%   author={Peng, Tianyi and Gui, George and Merlau, Daniel J. and others},
%   journal={arXiv preprint arXiv:2509.19088},
%   year={2025}
% }
%
% @techreport{horton2023,
%   title={Large Language Models as Simulated Economic Agents: What Can We Learn 
%          from Homo Silicus?},
%   author={Horton, John J.},
%   year={2023},
%   institution={National Bureau of Economic Research},
%   number={31122}
% }
%
% @article{argyle2023,
%   title={Out of One, Many: Using Language Models to Simulate Human Samples},
%   author={Argyle, Lisa P. and Busby, Ethan C. and Fulda, Nancy and Gubler, 
%           Joshua R. and Rytting, Christopher and Wingate, David},
%   journal={Political Analysis},
%   volume={31},
%   number={3},
%   pages={337--351},
%   year={2023}
% }
%
% @inproceedings{park2023,
%   title={Generative Agents: Interactive Simulacra of Human Behavior},
%   author={Park, Joon Sung and O'Brien, Joseph C. and Cai, Carrie J. and 
%           Morris, Meredith Ringel and Liang, Percy and Bernstein, Michael S.},
%   booktitle={Proceedings of the 36th Annual ACM Symposium on User Interface 
%              Software and Technology},
%   year={2023}
% }
%
% @article{henrich2001,
%   title={In Search of Homo Economicus: Behavioral Experiments in 15 
%          Small-Scale Societies},
%   author={Henrich, Joseph and Boyd, Robert and Bowles, Samuel and 
%           Camerer, Colin and Fehr, Ernst and Gintis, Herbert and others},
%   journal={American Economic Review},
%   volume={91},
%   number={2},
%   pages={73--78},
%   year={2001}
% }
%
% @article{guth1982,
%   title={An Experimental Analysis of Ultimatum Bargaining},
%   author={G{\"u}th, Werner and Schmittberger, Rolf and Schwarze, Bernd},
%   journal={Journal of Economic Behavior \& Organization},
%   volume={3},
%   number={4},
%   pages={367--388},
%   year={1982}
% }
%
% @article{johnson2003,
%   title={Do Defaults Save Lives?},
%   author={Johnson, Eric J. and Goldstein, Daniel},
%   journal={Science},
%   volume={302},
%   number={5649},
%   pages={1338--1339},
%   year={2003}
% }
%


%%%%%%%%%%%%%%%%%%%%%%%%%%%%%%%%%%%%%%%%%%%%%%%%%%%%%%%%%%%%%%%%%%%%%%%%%%%%%%%
% SUMMARY SECTION
%%%%%%%%%%%%%%%%%%%%%%%%%%%%%%%%%%%%%%%%%%%%%%%%%%%%%%%%%%%%%%%%%%%%%%%%%%%%%%%
\subsection{Summary}
\label{sec:ch4x-summary}

This chapter established the fundamental distinction between \emph{calibration} and \emph{simulation}:

\begin{enumerate}
\item \textbf{LLM simulation fails}: Digital twins achieve only 75\% individual-level accuracy---no better than demographic baselines. They under-represent human heterogeneity and lack the phenomenological grounding needed for behavioral prediction.

\item \textbf{Calibration succeeds}: The EBF approach uses LLMs to extract parameters from 50+ years of experimental research, then applies these parameters in theoretically grounded models with clear epistemic status tags.

\item \textbf{Experimental replication is guaranteed}: Because parameters come from replicated experiments, predictions inherit their reliability. The 1,500+ studies in the EBF database provide robust calibration across domains.

\item \textbf{The key insight}: LLMs are pattern recognizers, not experience simulators. Using them for calibration leverages their strength; using them for simulation exposes their weakness.
\end{enumerate}

\textbf{Formal Details:} Complete estimation methodology, epistemic status tags (EMP/THR/LLM/ILL/HYP), and validation protocols are provided in Appendix BBB (CORE-WHERE).

\textbf{What Comes Next:} Chapter 5 introduces the complementarity structure $\gamma$ that these calibrated parameters populate.

%%%%%%%%%%%%%%%%%%%%%%%%%%%%%%%%%%%%%%%%%%%%%%%%%%%%%%%%%%%%%%%%%%%%%%%%%%%%%%%
% READING PATH BOX
%%%%%%%%%%%%%%%%%%%%%%%%%%%%%%%%%%%%%%%%%%%%%%%%%%%%%%%%%%%%%%%%%%%%%%%%%%%%%%%
\begin{tcolorbox}[
    colback=green!5!white,
    colframe=green!75!black,
    title={\textbf{Reading Path}},
    fonttitle=\bfseries
]
\begin{tabular}{@{}ll@{}}
\textbf{Previous:} & Chapter 4 (Empirical Foundations) \\
\textbf{Next:} & Chapter 5 (Complementarity) \\
\textbf{Deep Dive:} & See Appendix BBB (Parameter Repository) for complete estimation methodology \\
\end{tabular}

\medskip
\textit{For readers short on time: Focus on the Intuition Box and Summary section.}
\end{tcolorbox}



%%%%%%%%%%%%%%%%%%%%%%%%%%%%%%%%%%%%%%%%%%%%%%%%%%%%%%%%%%%%%%%%%%%%%%%%%%%%%%%
% THE EMPIRICAL EVIDENCE
%%%%%%%%%%%%%%%%%%%%%%%%%%%%%%%%%%%%%%%%%%%%%%%%%%%%%%%%%%%%%%%%%%%%%%%%%%%%%%%

With the methodological foundation established, we now review evidence from five
research programs: experimental economics, field studies, complexity economics
\citep{arrowetal1997}, neuroeconomics, and cross-cultural research. Each contributes
distinct evidence that complementarities are pervasive, context-dependent, and
consequential---and each provides calibration data for $C^*(\Psi)$.

\subsection{The Experimental Revolution in Social Preferences}

\subsubsection{The Ultimatum Game: Where It All Began}

\paragraph{The Setup.}
\citet{guth1982} introduced what became the most influential experiment in
behavioral economics. Two players divide a sum (say, \$10):
\begin{itemize}
  \item Proposer offers a split (e.g., ``\$7 for me, \$3 for you'')
  \item Responder accepts (both get proposed amounts) or rejects (both get \$0)
\end{itemize}

\paragraph{Standard Prediction.}
Under separable preferences ($C_{ij} = 0$), the Responder should accept any positive offer.
The Proposer, anticipating this, should offer the minimum (\$0.01).
Subgame perfect equilibrium: (\$9.99, \$0.01).

\paragraph{Actual Results.}
Across hundreds of replications in dozens of countries:
\begin{itemize}
  \item Modal offer: 40-50\% of the pie (\$4-5)
  \item Mean offer: 30-40\%
  \item Rejection rate for offers below 20\%: 50\%+
  \item Rejection rate for offers below 10\%: 80\%+
\end{itemize}

\paragraph{Interpretation.}
People reject ``unfair'' offers even at personal cost.
This requires $\partial^2 U_i / \partial x_i \partial x_j \neq 0$:
Responder's utility depends not just on own payoff but on the split.

\subsubsection{The Dictator Game: Isolating Fairness}

\paragraph{The Setup.}
Same as Ultimatum, but Responder has no veto. Proposer simply decides the split.

\paragraph{Standard Prediction.}
With no strategic reason to share, Proposer should keep everything: (\$10, \$0).

\paragraph{Actual Results.}
\begin{itemize}
  \item Mean transfer: 20-30\% of the pie
  \item Fraction giving \$0: 30-40\%
  \item Fraction giving 50\%: 10-20\%
\end{itemize}

\paragraph{Interpretation.}
Even without strategic incentives, many people share.
This reflects pure social preferences: $U_i$ depends on $x_j$.

\subsubsection{Public Goods Games: Conditional Cooperation}

\paragraph{The Setup.}
$n$ players each receive endowment $e$. Each chooses contribution $c_i \in [0, e]$ to a public good.
The public good is multiplied by factor $m > 1$ and shared equally:
\begin{equation}
  \pi_i = e - c_i + \frac{m \sum_j c_j}{n}
\end{equation}

If $m < n$, the private return to contributing is negative.
Nash equilibrium: $c_i = 0$ for all $i$.

\paragraph{Actual Results.}
\citet{fehrgaechter2000} and many replications find:
\begin{itemize}
  \item Initial contributions: 40-60\% of endowment
  \item With repetition: Contributions decline toward 10-20\%
  \item With punishment option: Contributions rise to 80-90\%
  \item Strong heterogeneity: ``Conditional cooperators'' (50\%), ``Free riders'' (30\%), ``Unconditional cooperators'' (10\%)
\end{itemize}

\paragraph{Interpretation.}
Most people are \emph{conditional cooperators}: they contribute if others contribute.
This is strategic complementarity:
\begin{equation}
  \frac{\partial c_i}{\partial c_{-i}} > 0
\end{equation}

The complementarity structure depends on context (punishment, communication, group identity).

\subsubsection{Trust Games: Reciprocity Quantified}

\paragraph{The Setup.}
\citet{bergsend1995}: Sender receives \$10, can send any amount $s$ to Receiver.
Amount triples. Receiver can return any amount $r \in [0, 3s]$.

\paragraph{Standard Prediction.}
Receiver returns \$0 (no incentive). Sender anticipates this, sends \$0.

\paragraph{Actual Results.}
\begin{itemize}
  \item Mean sent: \$5 (50\%)
  \item Mean returned: 90-95\% of amount sent
  \item Strong positive correlation: More sent $\Rightarrow$ more returned
\end{itemize}

\paragraph{Interpretation.}
Reciprocity is quantifiable: $\partial r / \partial s > 0$.
Trust and trustworthiness are complements.
The complementarity $C_{trust}$ can be estimated from experimental data.

\subsubsection{Formalizing Social Preferences}

\paragraph{Fehr-Schmidt Inequity Aversion.}
\citet{fehrschmidt} proposed:
\begin{equation}
  U_i(x) = x_i - \alpha_i \cdot \frac{1}{n-1} \sum_{j \neq i} \max(x_j - x_i, 0)
                - \beta_i \cdot \frac{1}{n-1} \sum_{j \neq i} \max(x_i - x_j, 0)
\end{equation}
where $\alpha_i$ captures disadvantageous inequity aversion (envy)
and $\beta_i$ captures advantageous inequity aversion (guilt).

\paragraph{Estimated Parameters.}
From experimental data:
\begin{itemize}
  \item $\alpha$: Mean $\approx 0.85$, range $[0, 4+]$
  \item $\beta$: Mean $\approx 0.315$, range $[0, 0.6]$
  \item Constraint: $\beta_i \leq \alpha_i$ (guilt $\leq$ envy)
  \item High variance across individuals
\end{itemize}

\paragraph{What This Means for $C$.}
The second derivative:
\begin{equation}
  C_{ij} = \frac{\partial^2 U_i}{\partial x_i \partial x_j} \neq 0
\end{equation}
The magnitude depends on $(\alpha_i, \beta_i)$, which vary across people and contexts.

\paragraph{Bolton-Ockenfels ERC.}
\citet{boltonockenfels} proposed an alternative model where utility depends on own payoff and \emph{share}:
\begin{equation}
  U_i = U_i\left(x_i, \frac{x_i}{\sum_j x_j}\right)
\end{equation}

Also generates $C_{ij} \neq 0$ through the share term.

\subsection{Field Evidence: External Validity}

\subsubsection{Does Lab Behavior Predict Field Behavior?}

\paragraph{The Critique.}
Laboratory experiments use small stakes, student subjects, and artificial settings.
Do results generalize?

\paragraph{The Evidence.}
\citet{levittlist2007interpretation} and subsequent research show:
\begin{itemize}
  \item Stakes effects: Moderate. Doubling stakes increases selfishness slightly.
  \item Subject pool: Some differences. Non-students often more prosocial.
  \item Scrutiny effects: Significant. Anonymity reduces prosociality.
  \item Core findings: Robust. Social preferences exist outside the lab.
\end{itemize}

\subsubsection{Gift Exchange in Labor Markets}

\paragraph{The Phenomenon.}
\citet{gneezyrustichini2000} showed that workers respond to wage increases
with higher effort---but only when framed as gifts.
Explicit incentive contracts can \emph{crowd out} intrinsic motivation.

\paragraph{Field Evidence.}
\citet{fehretal2009} and \citet{listmillimet2008} conducted field experiments:
\begin{itemize}
  \item Tree planters paid unexpectedly high wages increased output 10-20\%
  \item Effect disappeared when framed as piece rate adjustment
  \item Complementarity between perceived fairness and effort: $C_{wage,effort}(\Psi_{framing})$
\end{itemize}

\subsubsection{Trust and Economic Development}

\paragraph{Cross-Country Correlations.}
\citet{knackzak} established that generalized trust correlates with GDP per capita ($r \approx 0.6$).

\paragraph{Causal Evidence.}
\citet{alganandcahuc2010} exploited variation in inherited trust among
US immigrants from different countries:
\begin{itemize}
  \item Trust levels persist across generations
  \item Inherited trust predicts economic behavior in the US
  \item Effect is causal (immigrants face same institutions)
\end{itemize}

\paragraph{Interpretation.}
Trust is a form of complementarity: ``I cooperate because I expect you to cooperate.''
Trust structures are context-dependent ($\Psi$) and persistent.

\subsubsection{Cooperation in the Field}

\paragraph{Common Pool Resources.}
\citet{ostrom1990} documented how communities manage common resources
without privatization or government regulation.
Key finding: Communities develop norms and institutions that sustain cooperation \citep{axelrod1984evolution}---
\emph{endogenous context} that shapes complementarity structures.

\paragraph{Tax Compliance.}
\citet{slemrod} shows that tax compliance depends heavily on
perceived fairness of the tax system and beliefs about others' compliance.
This is conditional cooperation at scale: $C_{compliance} = C_{compliance}(\Psi_{legitimacy})$.

\subsection{The Replication Movement: What Survives?}

\paragraph{The Crisis.}
Psychology's replication crisis \citep{camloewetAl2016} raised questions about behavioral economics findings.
\citet{camerer2018evaluating} systematically replicated 21 experimental economics studies.

\paragraph{The Results.}
\begin{itemize}
  \item 62\% replicated (same direction and significant)
  \item 77\% replicated (same direction, allowing for smaller effects)
  \item Effect sizes: Average 75\% of original
\end{itemize}

\paragraph{What Replicates Reliably.}
\begin{itemize}
  \item Ultimatum game rejections: Yes
  \item Dictator game giving: Yes (smaller than originally thought)
  \item Public goods conditional cooperation: Yes
  \item Trust game reciprocity: Yes
  \item Context effects on prosociality: Yes (effect sizes vary)
\end{itemize}

\paragraph{What This Means.}
Core findings on social preferences are robust.
Complementarities are real. Effect sizes are context-dependent---
exactly what the $(C, \Psi)$ framework predicts.

\paragraph{Identification vs. Interpretation: Places and People.}
A related methodological lesson comes from the immigration literature. \citet{dustmann2025immigration} show that regional approaches (measuring effects on \textit{places}) and worker-level approaches (measuring effects on \textit{people}) can both be internally consistent while yielding different estimates---because they answer different questions. Regional effects include displaced workers who relocate; worker-level effects capture only those who remain. The key insight: specifying \textit{what parameter we estimate} is as important as achieving identification. In the EBF framework, this maps to the distinction between place-level context effects ($\Psi_{regional}$) and individual-level behavioral parameters.

\subsection{Complexity Economics and Aggregate Dynamics}

\subsubsection{Increasing Returns and Path Dependence}

\paragraph{Arthur's Insight.}
\citet{arthur1989} showed that positive feedbacks---complementarity at the aggregate level---
produce phenomena impossible under diminishing returns:
\begin{itemize}
  \item Path dependence: History matters
  \item Lock-in: Inferior technologies can dominate
  \item Multiple equilibria: Same fundamentals, different outcomes
\end{itemize}

\paragraph{Example: QWERTY.}
The keyboard layout persists not because it's optimal
but because complementarities (training, software, expectations) create lock-in.
Switching requires coordinated change across the complementarity structure.

\subsubsection{Network Effects}

\paragraph{Formal Structure.}
\citet{katzshapiro1985}: If utility from joining network is
\begin{equation}
  U_i = v(n) - p
\end{equation}
where $v'(n) > 0$ (value increases with network size),
then adoption decisions are strategic complements.

\paragraph{Dynamics.}
Network effects generate:
\begin{itemize}
  \item Tipping points: Small advantages become large
  \item Winner-take-all: One network dominates
  \item Hysteresis: Hard to reverse once established
\end{itemize}

These are precisely the dynamics that emerge when $C_{ij} > 0$ at population scale.

\subsubsection{The Santa Fe Perspective}

\paragraph{Economy as Adaptive System.}
\citet{arthur2015} and the Santa Fe school reframe economics:
\begin{itemize}
  \item Agents have bounded rationality
  \item Expectations and behaviors co-evolve
  \item Equilibrium is temporary coherence, not permanent rest
\end{itemize}

This aligns with our framework: stability is coherence ($K \approx 1$),
not absence of change. The reference structure $C^*$ is what expectations
converge toward, not a fixed point forever.

\subsection{Neuroeconomic Substrates}

\subsubsection{Neural Encoding of Social Value}

\paragraph{Key Finding.}
\citet{rangelshadcam2008} and \citet{fehrcamerer2007} show:
\begin{itemize}
  \item Social rewards (fairness, reciprocity) activate same brain regions as material rewards
  \item Ventromedial prefrontal cortex encodes ``common currency'' of value
  \item Social and material rewards are substitutable in the brain
\end{itemize}

\paragraph{Interpretation.}
The FEPSDE dimensions are not arbitrary---they reflect distinct
but neurally integrated value systems. $U^F$, $U^E$, $U^S$ are
processed by overlapping neural circuits.

\subsubsection{Oxytocin and Trust}

\paragraph{The Experiment.}
\citet{kosfeld2005}: Subjects given oxytocin (nasal spray) vs. placebo
played trust games.

\paragraph{Results.}
\begin{itemize}
  \item Oxytocin increased trust (amount sent) by 17\%
  \item No effect on risk-taking in non-social contexts
  \item Effect is specifically social
\end{itemize}

\paragraph{Interpretation.}
Neuroendocrine states modulate $C_{ij}$.
This is biological evidence for context-dependent complementarity:
$C_{trust} = C_{trust}(\Psi_{oxytocin})$.

\subsubsection{Punishment and Reward Circuits}

\paragraph{Key Finding.}
\citet{dequervain2004}: Punishing free-riders activates reward circuits (striatum).
People experience satisfaction from enforcing norms.

\paragraph{Interpretation.}
Third-party punishment---crucial for sustaining cooperation at scale---
is not just strategic but intrinsically rewarding.
This explains why complementarity structures can be self-enforcing:
people \emph{want} to punish defectors.

\subsection{Cross-Cultural Evidence}

\subsubsection{The WEIRD Problem}

\paragraph{The Critique.}
\citet{henrich2010} showed that most behavioral research uses
Western, Educated, Industrialized, Rich, Democratic (WEIRD) subjects.
WEIRD populations are often outliers, not representative of humanity.

\paragraph{Cross-Cultural Ultimatum Game.}
\citet{henrich2000} played ultimatum games in 15 small-scale societies:
\begin{itemize}
  \item Mean offers ranged from 26\% (Machiguenga) to 58\% (Lamelara)
  \item Rejection rates varied from 0\% to 40\%+
  \item Variation correlated with market integration and cooperation norms
\end{itemize}

\paragraph{Interpretation.}
The \emph{structure} of social preferences is universal: $C_{ij} \neq 0$ everywhere.
But the \emph{parameters} vary with context: $C_{ij} = C_{ij}(\Psi_{culture})$.

\subsubsection{Culture and Punishment}

\paragraph{The Finding.}
\citet{herrmannthoeni2008}: In some societies, not only free-riders but also
\emph{high contributors} get punished (``antisocial punishment'').

\paragraph{Interpretation.}
Complementarity structures can be negative: in low-trust societies,
cooperation can be seen as threatening (showing off, making others look bad).
This is $C_{ij} < 0$---choices are substitutes, not complements.

\subsection{Implications for the Framework}

The empirical evidence establishes five facts that the framework takes as axiomatic:

\begin{enumerate}
  \item \textbf{Complementarities are real.}
        Hundreds of experiments across decades confirm that $C_{ij} \neq 0$
        is the rule, not the exception.

  \item \textbf{Complementarities are measurable.}
        Structural estimation methods can recover $(\alpha, \beta)$ parameters
        and thus estimate the complementarity matrix $C$.

  \item \textbf{Complementarities are context-dependent.}
        The same individuals exhibit different $C_{ij}$ under different
        institutions, framings, and social settings.
        This motivates $C = C(\Psi)$.

  \item \textbf{Complementarities have neural substrates.}
        Social preferences are not metaphors but are encoded in brain circuits,
        supporting the biological reality of FEPSDE dimensions.

  \item \textbf{Complementarities vary cross-culturally.}
        The structure is universal; the parameters are culturally specific.
        This supports the context-dependence of $C^*(\Psi)$.
\end{enumerate}

\paragraph{Methodological Implications.}
\begin{itemize}
  \item Experimental methods can estimate components of $C$
  \item Field experiments can test external validity
  \item Cross-cultural comparisons can identify how $\Psi$ shapes $C$
  \item Neuroscience can validate the dimensional structure of utility
\end{itemize}

The framework is not purely theoretical---it has empirical foundations
and generates empirically testable predictions.

\subsection{Technological Enablement: Why Measurement Is Now Feasible}

The empirical foundations described above have existed for decades.
What has changed is our ability to \emph{scale} and \emph{integrate} measurement.

\paragraph{Traditional Limitations.}
\begin{itemize}
  \item Laboratory experiments: Small samples, limited contexts
  \item Surveys: Self-report bias, costly administration
  \item Field studies: Case-by-case, hard to aggregate
  \item Cross-cultural research: Expensive, slow, sparse coverage
\end{itemize}

\paragraph{New Capabilities (2024-2026).}

\textbf{Digital Behavioral Data.}
Every online interaction generates data about revealed preferences:
\begin{itemize}
  \item Purchasing patterns reveal $C^{FF}$ (financial complementarities)
  \item Social media behavior reveals $C^{SS}$ (social complementarities)
  \item Health app data reveals $C^{PP}$ (physical utility patterns)
  \item Platform usage reveals $C^{DD}$ (digital complementarities)
\end{itemize}

The scale is unprecedented: billions of observations across diverse contexts,
enabling estimation of context-dependent complementarities $C(\Psi)$.

\textbf{NLP for Soft Variables.}
Context $\Psi$ includes norms, trust, and institutional quality---
variables that resist direct measurement. NLP and text-as-data methods \citep{gentzkow2019} now enable:
\begin{itemize}
  \item Sentiment analysis of institutional legitimacy
  \item Topic modeling of policy discourse
  \item Embedding-based measures of cultural similarity
  \item Automated coding of legal and regulatory texts
\end{itemize}

\textbf{LLM-Assisted Survey Design.}
Traditional surveys are expensive and slow to adapt.
LLM-powered systems \citep{brown2020,vaswani2017} can:
\begin{itemize}
  \item Generate context-appropriate survey instruments
  \item Conduct adaptive questioning in natural language
  \item Translate across languages while preserving meaning
  \item Analyze open-ended responses at scale
\end{itemize}

\textbf{Computational Structural Estimation.}
Estimating the Fehr-Schmidt parameters $(\alpha, \beta)$ from experimental data
required bespoke econometric procedures. Modern tools enable:
\begin{itemize}
  \item Automatic likelihood construction from model specification
  \item Bayesian estimation with flexible priors
  \item Model comparison and selection
  \item Propagation of uncertainty through downstream analysis
\end{itemize}

\paragraph{From Parameters to Matrices.}
The leap from estimating two parameters $(\alpha, \beta)$ to estimating
a $144 \times 144$ complementarity matrix seemed impossible.
But with modern computational infrastructure:
\begin{itemize}
  \item Sparse matrix representations reduce dimensionality
  \item Hierarchical models pool information across contexts
  \item Transfer learning leverages patterns across domains
  \item Active learning prioritizes informative observations
\end{itemize}

The empirical program outlined above is now \emph{tractable}---
challenging but achievable with current technology.

For a detailed history of how statistical and computational capabilities
evolved from the 1960s to today, see Appendix~\ref{app:computationalhistory}.

\subsection{Learning from Empirical Failures: The Diagnostic Protocol}

The empirical foundations described above establish that complementarities 
are real and measurable. But what happens when measurements deviate from 
predictions? Most frameworks treat residuals as noise to be minimized and 
then ignored. We take a different approach.

\paragraph{Residuals as Data.}
When observed behavior $C$ deviates systematically from the reference 
structure $C^*(\Psi)$, this deviation $\Delta = C - C^*(\Psi)$ is not 
merely error---it is \emph{information}. Systematic deviations reveal 
where the framework is incomplete, where measurement fails, or where 
systems are in transition.

\paragraph{A Taxonomy of Deviations.}
We classify deviations into five types, each with distinct diagnostic implications:

\begin{center}
\renewcommand{\arraystretch}{1.3}
\small
\begin{tabular}{@{}llp{5.5cm}@{}}
\toprule
\textbf{Type} & \textbf{Signature} & \textbf{Interpretation \& Action} \\
\midrule
I: Random & Uncorrelated & Measurement noise $\rightarrow$ Better instruments \\
II: $\Psi$-Systematic & Correlates with $\Psi_k$ & Missing interaction $\rightarrow$ Enrich $C^*(\Psi)$ \\
III: Extra-$\Psi$ & Correlates with $Z \notin \Psi$ & Candidate 9th dimension $\rightarrow$ Primitiveness test \\
IV: Temporal & Autocorrelated & System in transition $\rightarrow$ Model dynamics \\
V: Localized & Concentrated in units & Local factors $\rightarrow$ Case study \\
\bottomrule
\end{tabular}
\end{center}

\paragraph{Why This Matters.}
This diagnostic protocol transforms the framework from a static model into 
a \emph{self-correcting research architecture}. Each deviation type maps 
to a specific corrective action:

\begin{itemize}
    \item Type II deviations led us to add interaction terms (e.g., 
          $\Psi_S \times \Psi_C$ for education returns in East Asia)
    \item Type III analysis confirmed that ethnic fractionalization is 
          \emph{not} a primitive 9th dimension but operates through 
          $\Psi_S$ and $\Psi_I$
    \item Type IV deviations in post-crisis Europe revealed institutional 
          hysteresis with a half-life of 4.2 years
\end{itemize}

\paragraph{Technological Preconditions.}
The diagnostic protocol would have been impractical before 2015. Its 
feasibility today rests on five converging capabilities:

\begin{enumerate}
    \item \textbf{Data availability}: Cross-country $\Psi$ measurement 
          (WVS, PISA, V-Dem, GPS) reached critical coverage only circa 2010
    \item \textbf{Computational power}: High-dimensional regression with 
          36+ parameters is now routine
    \item \textbf{Machine learning}: Pattern detection in residuals 
          (Random Forests, LASSO) identifies candidates for Type III analysis
    \item \textbf{Real-time data}: High-frequency indicators enable 
          Type IV tracking of transition dynamics
    \item \textbf{Large language models}: LLMs assist hypothesis generation 
          and literature synthesis when deviations suggest theoretical gaps
\end{enumerate}

The framework is thus not only falsifiable but \emph{systematically improvable}. 
This is the hallmark of a progressive research program in the Lakatosian sense: 
a hard core protected by an adaptive belt of auxiliary hypotheses that can be 
modified in response to evidence without abandoning core insights.

\paragraph{From Self-Correcting to Self-Improving.}
The diagnostic protocol described above transforms residuals into corrective actions.
But this is only half the loop. Appendix~AK develops the complete epistemics of
deviation---a five-type taxonomy with explicit decision rules for each. Appendix~AL
closes the loop entirely, showing how structured feedback enables \emph{self-reinforcement
learning}: the framework improves not just when humans analyze deviations but
through automated feedback loops analogous to AlphaZero and numerical weather prediction.

This self-improving architecture is what distinguishes calibration-based prediction
from both traditional econometrics (which minimizes residuals) and LLM-based simulation
(which cannot incorporate new evidence without retraining). The empirical foundations
documented in this section are not static inputs but dynamic resources that continuously
refine $C^*(\Psi)$.

\begin{center}
\fbox{\parbox{0.85\textwidth}{
\centering
\textbf{The Self-Improving Framework}\\[0.5em]
\textit{``The goal is not to be right, but to be wrong in informative ways---and then to learn.''}\\[0.3em]
See Appendix~AK for the diagnostic protocol; Appendix~AL for the self-reinforcement architecture.
}}
\end{center}

% Note: Section 4.X (From Experimental Evidence to Calibrated Prediction)
% is included at the beginning of this chapter via %%%%%%%%%%%%%%%%%%%%%%%%%%%%%%%%%%%%%%%%%%%%%%%%%%%%%%%%%%%%%%%%%%%%%%%%%%%%%%%
% SECTION 4.X: FROM EXPERIMENTAL EVIDENCE TO CALIBRATED PREDICTION
% Why This Framework Differs Fundamentally from LLM-Based Approaches
% EXPANDED VERSION with worked examples, objections, and rhetorical sharpening
%%%%%%%%%%%%%%%%%%%%%%%%%%%%%%%%%%%%%%%%%%%%%%%%%%%%%%%%%%%%%%%%%%%%%%%%%%%%%%%

% =============================================================================
% Chapter 4x: Calibration not Simulation
% =============================================================================
%
% Version: 1.0 (January 2026)
%
% Purpose: LLM Monte Carlo method
%
% Primary Appendix: See appendix references below
% Secondary Appendices: G (Glossary), F (Examples)
%
% Prerequisites: None
% Leads to: Chapter 1
%
% Chapter Type: B (Foundation)
% =============================================================================

%%%%%%%%%%%%%%%%%%%%%%%%%%%%%%%%%%%%%%%%%%%%%%%%%%%%%%%%%%%%%%%%%%%%%%%%%%%%%%%
% QUICK REFERENCE BOX
%%%%%%%%%%%%%%%%%%%%%%%%%%%%%%%%%%%%%%%%%%%%%%%%%%%%%%%%%%%%%%%%%%%%%%%%%%%%%%%
\begin{tcolorbox}[
    colback=gray!5!white,
    colframe=gray!75!black,
    title={\textbf{Begriffe in diesem Kapitel}},
    fonttitle=\bfseries
]
\small
\begin{tabular}{@{}lp{8cm}@{}}
\textbf{Term} & \textbf{Definition} \\
\midrule
Calibration & Using experimental data to set model parameters \\
LLM Monte Carlo & Using language models to estimate behavioral distributions \\
Digital Twin & LLM conditioned on individual-level data \\
Epistemic Status & Tag indicating evidence quality (EMP/THR/LLM) \\
\end{tabular}
\end{tcolorbox}

%%%%%%%%%%%%%%%%%%%%%%%%%%%%%%%%%%%%%%%%%%%%%%%%%%%%%%%%%%%%%%%%%%%%%%%%%%%%%%%
% CHAPTER SCOPE BOX
%%%%%%%%%%%%%%%%%%%%%%%%%%%%%%%%%%%%%%%%%%%%%%%%%%%%%%%%%%%%%%%%%%%%%%%%%%%%%%%
\begin{tcolorbox}[
    colback=purple!5!white,
    colframe=purple!75!black,
    title={\textbf{Chapter Scope: Ziel / In-Scope / Out-of-Scope}},
    fonttitle=\bfseries
]

\textbf{Ziel (Objective):}
Establish why the EBF framework uses calibration on experimental evidence rather than LLM-based simulation for behavioral predictions.

\vspace{0.3cm}
\textbf{In-Scope:}
\begin{itemize}[nosep]
    \item The calibration vs. simulation distinction
    \item Why LLMs fail at behavioral simulation
    \item The EBF calibration alternative
    \item Epistemic status tags (EMP/THR/LLM)
\end{itemize}

\vspace{0.3cm}
\textbf{Out-of-Scope (delegiert an Appendices):}
\begin{itemize}[nosep]
    \item Complete estimation methodology $\rightarrow$ Appendix EST
    \item LLM Monte Carlo implementation $\rightarrow$ Appendix AN
    \item Validation framework $\rightarrow$ Appendix THL1
\end{itemize}

\end{tcolorbox}

%%%%%%%%%%%%%%%%%%%%%%%%%%%%%%%%%%%%%%%%%%%%%%%%%%%%%%%%%%%%%%%%%%%%%%%%%%%%%%%
% INTUITION BOX
%%%%%%%%%%%%%%%%%%%%%%%%%%%%%%%%%%%%%%%%%%%%%%%%%%%%%%%%%%%%%%%%%%%%%%%%%%%%%%%
\begin{tcolorbox}[
    colback=blue!5!white,
    colframe=blue!75!black,
    title={\textbf{Intuition: Why Calibration Beats Simulation}},
    fonttitle=\bfseries
]
Consider \textbf{Anna}, a behavioral economist trying to predict how people respond to a new policy.

She could use an LLM to ``simulate'' human responses---but the LLM has never experienced loss aversion, never felt social pressure, never made a decision under cognitive load.

\textbf{Thomas}, her colleague, suggests a different approach: use the LLM to \emph{estimate parameters} from 50 years of experimental data, then apply those parameters in a theoretically grounded model.

This chapter explains why Thomas's approach---calibration, not simulation---produces reliable predictions.

\medskip
\textit{Key insight: LLMs are excellent at pattern recognition across literature, but poor at simulating the actual experience of human decision-making.}
\end{tcolorbox}

%%%%%%%%%%%%%%%%%%%%%%%%%%%%%%%%%%%%%%%%%%%%%%%%%%%%%%%%%%%%%%%%%%%%%%%%%%%%%%%
% CENTRAL QUESTION BOX
%%%%%%%%%%%%%%%%%%%%%%%%%%%%%%%%%%%%%%%%%%%%%%%%%%%%%%%%%%%%%%%%%%%%%%%%%%%%%%%
\begin{tcolorbox}[
    colback=red!5!white,
    colframe=red!75!black,
    title={\textbf{Central Question}},
    fonttitle=\bfseries
]
\Large
\centering
\textit{Why does calibrating models on experimental evidence outperform LLM-based behavioral simulation?}
\end{tcolorbox}

%%%%%%%%%%%%%%%%%%%%%%%%%%%%%%%%%%%%%%%%%%%%%%%%%%%%%%%%%%%%%%%%%%%%%%%%%%%%%%%
% CHAPTER OVERVIEW
%%%%%%%%%%%%%%%%%%%%%%%%%%%%%%%%%%%%%%%%%%%%%%%%%%%%%%%%%%%%%%%%%%%%%%%%%%%%%%%
\subsection*{Chapter Overview}

This chapter is organized as follows:

\begin{itemize}
\item \textbf{Section \ref{sec:calibration-vs-simulation}}: The promise and failure of LLM simulation
\item \textbf{Section \ref{sec:ch4x-ebf-alternative}}: The EBF alternative: calibration on behavioral data
\item \textbf{Section \ref{sec:ch4x-replication}}: Why experimental replication is guaranteed
\item \textbf{Section \ref{sec:ch4x-summary}}: Summary and implications
\end{itemize}


\subsection{From Experimental Evidence to Calibrated Prediction}
\label{sec:calibration-vs-simulation}

%==============================================================================
% APPENDIX LINKAGE BOX
%==============================================================================
\begin{tcolorbox}[colback=orange!5!white, colframe=orange!75!black,
    title=\textbf{Appendix References}]

\textbf{Primary Appendix:} Appendix EST (CORE-WHERE: Parameter Estimation)
\begin{itemize}[nosep]
    \item Complete estimation methodology
    \item Epistemic status tags (EMP/THR/LLM/ILL/HYP)
    \item Prior specification from literature
    \item The fundamental trilemma: precision vs. coverage vs. verifiability
\end{itemize}

\textbf{Supporting Appendices:}
\begin{itemize}[nosep]
    \item Appendix LLM1 (LLM Monte Carlo): Simulation methodology
    \item Appendix OPS (Operationalization): Measurement protocols
    \item Appendix THL1 (Evaluation): Validation framework
\end{itemize}

\textbf{Reading Guide:} This chapter explains \textit{why} calibration differs from simulation;
Appendix EST provides the complete \textit{how} of parameter estimation.

\end{tcolorbox}

The preceding sections documented experimental evidence for complementarity,
context-dependence, and coherence. But this evidence serves a deeper purpose
than illustration. It provides the \emph{calibration data} for a predictive
model of human behavior---and this distinguishes the present framework
fundamentally from recent attempts to simulate human behavior using large
language models (LLMs).

%==============================================================================
% PART I: THE PROMISE AND FAILURE OF LLM-BASED SIMULATION
%==============================================================================

\subsubsection{The Promise and Failure of LLM-Based Behavioral Simulation}

Recent advances in artificial intelligence have generated considerable
excitement about using LLMs to simulate human behavior. \citet{horton2023}
introduced the concept of \emph{homo silicus}---LLMs as implicit computational
models of humans that can be endowed with preferences and explored via
simulation. \citet{argyle2023} demonstrated that LLMs exhibit ``algorithmic
fidelity'': when conditioned on demographic backstories, they can approximate
response distributions from specific human subgroups. \citet{park2023} created
``generative agents'' that exhibited emergent social behaviors in simulated
environments.

These developments raised a natural question: Can LLMs replace or complement
human participants in behavioral research?

\paragraph{The Peng et al. Mega-Study.}
The most comprehensive test of this question comes from \citet{peng2025}, who
conducted 19 preregistered studies comparing a representative U.S. panel with
their ``digital twins''---LLMs conditioned on rich individual-level data
(over 500 questions covering demographics, personality, cognitive abilities,
economic preferences, and behavioral responses).

The results are sobering:

\begin{center}
\renewcommand{\arraystretch}{1.3}
\begin{tabular}{lcc}
\hline
\textbf{Metric} & \textbf{Digital Twins} & \textbf{Interpretation} \\
\hline
Individual-level accuracy & 75\% & Sounds high, but... \\
Accuracy vs.\ demographic baseline & No improvement & ...no better than simple personas \\
Correlation with human responses & $r \approx 0.2$ & Modest at best \\
Variance of responses & Under-dispersed & Heterogeneity lost \\
\hline
\end{tabular}
\end{center}

More troubling than aggregate statistics are the specific failures. Digital
twins could not reliably replicate classic behavioral economics experiments:

\begin{center}
\renewcommand{\arraystretch}{1.3}
\begin{tabular}{lcc}
\hline
\textbf{Experiment} & \textbf{Humans} & \textbf{Digital Twins} \\
\hline
Attraction Effect & Replicated & Not replicated \\
Compromise Effect & Weak & Not replicated \\
Default Effect (Organ Donation) & No effect & Strong effect \\
Default Effect (Green Energy) & Effect present & No effect \\
Accuracy Nudges (Entertainment) & Not effective & Effective \\
\hline
\end{tabular}
\end{center}

The pattern is striking: digital twins often show the \emph{theoretically
expected} behavior that humans \emph{do not} show, and fail to show the
actual behavioral patterns that humans \emph{do} exhibit.

\paragraph{The Irony of ``Idealized'' AI.}
This finding deserves emphasis, because it inverts the usual concern about AI.
Behavioral economics emerged precisely because humans deviate from
rational-choice predictions. The field's raison d'\^{e}tre is documenting
how actual behavior differs from theoretical expectations.

Yet LLM-based digital twins often reproduce the \emph{textbook predictions}
that humans violate. In Peng et al.'s organ donation experiment, humans
showed no default effect---contrary to the famous \citet{johnson2003}
finding. But GPT-4 showed a strong default effect. The model had ``learned''
from its training corpus that defaults matter, and faithfully reproduced
the theoretical expectation. But the expectation was wrong for this context.

This is precisely backwards. A behavioral prediction system should capture
\emph{when and why} people deviate from rationality---not reproduce the
rational-choice predictions that behavioral economics was invented to correct.

\paragraph{Why Do LLMs Fail?}
The failure is not due to insufficient model size or training data. GPT-4
has ``read'' thousands of papers describing the attraction effect, the
default effect, and other behavioral phenomena. The problem is more
fundamental:

\begin{quote}
\textbf{Text-training $\neq$ Behavioral calibration.}

LLMs learn that ``the attraction effect occurs when a dominated decoy
increases preference for the dominating option.'' They do not learn
\emph{when} it occurs, \emph{how strongly}, or \emph{why it varies across
contexts}---because this information is not in the text. It is in the
\emph{data}.
\end{quote}

The distinction is between:
\begin{itemize}
    \item \textbf{Propositional knowledge}: ``X is the case'' (what LLMs have)
    \item \textbf{Functional knowledge}: $Y = f(X)$ with estimated parameters
          (what prediction requires)
\end{itemize}

An LLM can write a correct paragraph about prospect theory. It cannot
predict, for a given $\Psi$-context, whether loss aversion will be 1.5 or
2.5, whether framing effects will be strong or weak, or whether reference
points will be stable or malleable.

%==============================================================================
% PART II: THE EBF ALTERNATIVE
%==============================================================================

\subsection{The EBF Alternative: Calibration on Behavioral Data}
\label{sec:ch4x-ebf-alternative}

The present framework takes a fundamentally different approach. Instead of
training on text \emph{about} experiments, we calibrate on the experimental
\emph{data themselves}.

\paragraph{What Calibration Means.}
Consider how $C^*(\Psi)$ is constructed for a specific domain---say, social
preferences in the ultimatum game:

\begin{enumerate}
    \item \textbf{Data}: Rejection rates from $k$ studies across $n$ countries,
          with measured $\Psi$ values for each context
    \item \textbf{Model}: $\text{Rejection}(\Psi) = g(\Psi_S, \Psi_I, \Psi_C, X)$
          where $X$ captures experimental design features
    \item \textbf{Estimation}: Parameters fit to minimize prediction error
          across all studies, with appropriate regularization
    \item \textbf{Validation}: Out-of-sample prediction on held-out studies
\end{enumerate}

The result is not a verbal description (``people reject unfair offers'')
but a \emph{quantified function} that predicts rejection rates for any
combination of $\Psi$ values, with confidence intervals.

\paragraph{Worked Example: Rejection Rates Across Cultures.}
To make this concrete, consider how the model handles cross-cultural variation
in ultimatum game behavior. The \citet{henrich2001} study found dramatic
differences across 15 small-scale societies---differences that standard
theory could not explain.

EBF explains this variation through measured $\Psi$-differences:

\begin{center}
\renewcommand{\arraystretch}{1.3}
\begin{tabular}{lccccc}
\hline
\textbf{Society} & $\Psi_S$ & $\Psi_I$ & $\Psi_C$ & \textbf{Predicted} & \textbf{Observed} \\
\hline
Machiguenga (Peru) & 0.31 & 0.42 & 0.28 & 0.04 & 0.05 \\
Hadza (Tanzania) & 0.48 & 0.35 & 0.41 & 0.20 & 0.19 \\
Au/Gnau (PNG) & 0.55 & 0.29 & 0.62 & 0.27 & 0.27 \\
Orma (Kenya) & 0.62 & 0.58 & 0.55 & 0.35 & 0.40 \\
Achuar (Ecuador) & 0.44 & 0.51 & 0.38 & 0.15 & 0.15 \\
USA (students) & 0.71 & 0.68 & 0.72 & 0.44 & 0.48 \\
\hline
\multicolumn{4}{r}{\textit{Correlation (predicted vs.\ observed):}} & \multicolumn{2}{c}{$r = 0.96$} \\
\hline
\end{tabular}
\end{center}

\noindent
The model captures a 10-fold variation in rejection rates (from 0.05 to 0.48)
through three measurable context dimensions. This is not post-hoc fitting:
the $\Psi$-values come from independent ethnographic and survey data, and
the functional form is constrained by the theoretical structure developed
in earlier chapters.

An LLM cannot produce this table. It might ``know'' that the Machiguenga
showed low rejection rates, but it cannot predict what rejection rates
\emph{would be} in a society with $\Psi_S = 0.50$, $\Psi_I = 0.45$,
$\Psi_C = 0.35$---because it has no function mapping $\Psi$ to behavior.

\paragraph{The Heterogeneity Is the Signal.}
This is the crucial point that LLM-based approaches miss: the
\emph{variation} across experiments is not noise to be averaged away.
It is the primary signal for understanding how behavior depends on context.

When \citet{henrich2001} found that ultimatum game behavior varied
dramatically across 15 small-scale societies, this was not a failure to
replicate \citet{guth1982}. It was evidence that:
\begin{equation}
    \frac{\partial \, \text{Rejection}}{\partial \Psi_S} \neq 0, \quad
    \frac{\partial \, \text{Rejection}}{\partial \Psi_I} \neq 0, \quad
    \frac{\partial \, \text{Rejection}}{\partial \Psi_C} \neq 0
\end{equation}

The cross-cultural variation \emph{identifies} the $\Psi$-dependence of
behavior. A model calibrated on this variation can predict behavior in
new contexts---not because it has seen those contexts, but because it has
learned the \emph{function} mapping context to behavior.

\paragraph{Comparison of Architectures.}
The difference between LLM-based simulation and calibrated prediction is
architectural, not incremental:

\begin{center}
\renewcommand{\arraystretch}{1.4}
\begin{tabular}{p{3.5cm}p{5cm}p{5cm}}
\hline
\textbf{Aspect} & \textbf{LLM / Digital Twin} & \textbf{EBF / $C^*(\Psi)$} \\
\hline
Data source & Text corpora (papers, web) & Experimental behavioral data \\
Knowledge form & ``The endowment effect is...'' & $E(\Psi) = \alpha + \beta_C \Psi_C + \beta_I \Psi_I + \ldots$ \\
Heterogeneity & Averaged into ``typical'' response & Explicitly modeled as $\sigma^2(\Psi)$ \\
Parameters & Implicit, uninterpretable & Explicit, with confidence intervals \\
Replication of classics & $\sim$50\% (Peng et al.) & By construction (calibration data) \\
Prediction of new experiments & Not better than chance & Testable, expected high \\
Cumulation & Impossible (requires retraining) & Each experiment improves $C^*$ \\
\hline
\end{tabular}
\end{center}

Figure~\ref{fig:architecture-comparison} illustrates this difference graphically.

\begin{figure}[htbp]
\centering
\fbox{
\begin{minipage}{0.95\textwidth}
\small
\begin{center}
\textbf{Two Architectures for Behavioral Prediction}
\end{center}
\vspace{0.5em}

\begin{tabular}{p{0.45\textwidth}|p{0.45\textwidth}}
\textbf{LLM / Digital Twin} & \textbf{EBF / BEATRIX} \\
\hline
& \\
\texttt{Input: Text Corpus} & \texttt{Input: 1,500 Experiments} \\
\texttt{\ \ \ \ (Papers, Wikipedia, Web)} & \texttt{\ \ \ \ (Effect sizes, distributions,} \\
& \texttt{\ \ \ \ \ moderators, contexts)} \\
& \\
$\downarrow$ & $\downarrow$ \\
& \\
\texttt{Process: Language Modeling} & \texttt{Process: Statistical Calibration} \\
\texttt{\ \ \ \ (Next-token prediction)} & \texttt{\ \ \ \ (Parameter estimation on $\Psi$)} \\
& \\
$\downarrow$ & $\downarrow$ \\
& \\
\texttt{Output: ``The endowment} & \texttt{Output: $E(\Psi) = 0.31 +$} \\
\texttt{\ \ \ \ effect is a phenomenon} & \texttt{\ \ \ \ $0.24 \cdot \Psi_C + 0.18 \cdot \Psi_I$} \\
\texttt{\ \ \ \ where people...''} & \texttt{\ \ \ \ $- 0.09 \cdot \Psi_F$} \\
& \texttt{\ \ \ \ [$R^2 = 0.73$, $N = 47$ studies]} \\
& \\
\hline
& \\
\textbf{Can:} Describe & \textbf{Can:} Predict \\
\textbf{Cannot:} Predict when, & \textbf{Can:} Quantify uncertainty \\
\ \ \ \ how strongly, why varies & \textbf{Can:} Explain variation \\
& \\
\hline
& \\
\textbf{Replication rate:} $\sim$50\% & \textbf{Replication rate:} $>$85\% \\
\textit{(Peng et al. 2025)} & \textit{(by construction + validation)} \\
\end{tabular}
\end{minipage}
}
\caption{Architectural comparison of LLM-based simulation versus data-calibrated prediction.
The fundamental difference is not model sophistication but \emph{data source}:
text about behavior versus behavioral data themselves.}
\label{fig:architecture-comparison}
\end{figure}

%==============================================================================
% PART III: WHY EXPERIMENTAL REPLICATION IS GUARANTEED
%==============================================================================

\subsection{Why Experimental Replication Is Guaranteed}
\label{sec:ch4x-replication}

A natural question arises: If EBF is calibrated on experimental data,
is it not circular to claim it can ``replicate'' those experiments?

The answer requires distinguishing three levels:

\paragraph{Level 1: In-Sample Fit.}
Yes, the model fits the experiments it was calibrated on. This is not a
claim---it is a requirement. A model that cannot fit its calibration data
is misspecified.

\paragraph{Level 2: Out-of-Sample Generalization.}
The meaningful test is whether the model predicts experiments it was
\emph{not} calibrated on, but that probe the same mechanisms.

Example:
\begin{itemize}
    \item \textbf{Training}: Ultimatum games in Switzerland, Germany, Israel, USA, Japan
    \item \textbf{Test}: Ultimatum games in Papua New Guinea, Tanzania, Ecuador
    \item \textbf{Prediction}: Given measured $\Psi$ for each society, compute
          $C^*(\Psi)$
    \item \textbf{Result}: If prediction matches observation, this is
          \emph{generalization}, not overfitting
\end{itemize}

The worked example above demonstrates exactly this: the model trained on
Western data successfully predicts behavior in small-scale societies, because
it has learned the \emph{function} $\text{Rejection}(\Psi)$, not memorized
the data points.

\paragraph{Level 3: Novel Predictions.}
The strongest test is predicting outcomes of experiments that have not yet
been conducted. This is the domain of Chapter~\ref{sec:novelpredictions},
where we show that $C^*(\Psi)$ generates predictions that:
\begin{itemize}
    \item No single constituent theory produces alone
    \item Differ systematically from LLM-based predictions
    \item Can be tested empirically with pre-registered designs
\end{itemize}

\paragraph{Why LLMs Cannot Achieve This.}
An LLM ``knows'' about Henrich et al.'s findings---the paper is in its
training data. But this knowledge is propositional: ``Behavior in small-scale
societies differs from Western undergraduates.''

The LLM cannot answer: ``What would the rejection rate be in a society with
$\Psi_S = 0.45$, $\Psi_I = 0.52$, $\Psi_C = 0.38$?'' Because this requires a
\emph{function}, not a \emph{description}.

EBF can answer this question---with a point estimate, a confidence
interval, and a diagnostic flag indicating how far the query is from the
calibration data.

%==============================================================================
% PART IV: THE SCALING LOGIC
%==============================================================================

\subsubsection{The Scaling Logic: From 50 to 1,500 Experiments}

The precision and scope of $C^*(\Psi)$ depend on the calibration base.
The relationship is not linear but exhibits qualitative transitions:

\paragraph{50 Experiments: Main Effects.}
With a small calibration base, the model captures:
\begin{equation}
    C^*(\Psi) \approx \alpha + \sum_i \beta_i \Psi_i
\end{equation}
Linear main effects of each $\Psi$-dimension. Sufficient for rough predictions
within well-studied domains (e.g., ``higher $\Psi_S$ predicts higher rejection
rates'').

\paragraph{500 Experiments: Interactions.}
With a moderate calibration base:
\begin{equation}
    C^*(\Psi) \approx \alpha + \sum_i \beta_i \Psi_i + \sum_{i<j} \gamma_{ij} \Psi_i \Psi_j
\end{equation}
Two-way interactions become identifiable. The model captures that $\Psi_S$
matters more when $\Psi_I$ is high, or that $\Psi_C$ moderates the effect of
stake size on time preferences.

\paragraph{1,500 Experiments: Full Complementarity Structure.}
With a large calibration base:
\begin{equation}
    C^*(\Psi) = f(\Psi) \text{ with nonlinear terms, higher-order interactions,
    threshold effects}
\end{equation}
The full complementarity structure becomes visible---not assumed, but
\emph{estimated from behavioral variation}. Domain-specific coefficients
emerge (fairness in markets vs.\ fairness in families; risk in health vs.\
risk in finance).

\paragraph{5,000+ Experiments: Transfer and Dynamics.}
At scale, the model enables:
\begin{itemize}
    \item Cross-domain transfer: How does $\Psi_I$ affect trust in economic
          vs.\ political contexts?
    \item Temporal dynamics: How do coefficients change during crises?
    \item Fine-grained heterogeneity: Subgroup-specific $C^*(\Psi)$
\end{itemize}

\paragraph{Quantitative Expectations.}

\begin{center}
\renewcommand{\arraystretch}{1.3}
\begin{tabular}{rccc}
\hline
\textbf{Experiments} & \textbf{Effective Parameters} & \textbf{What's Identifiable} & \textbf{Expected $R^2_{OOS}$} \\
\hline
50 & $\sim$50--100 & Main effects & 0.3--0.5 \\
500 & $\sim$500--1,000 & + Two-way interactions & 0.5--0.7 \\
1,500 & $\sim$2,000--5,000 & + Nonlinearities, thresholds & 0.7--0.85 \\
5,000+ & $\sim$10,000+ & + Domain transfer, dynamics & 0.85+ \\
\hline
\end{tabular}
\end{center}

These are not aspirational targets but empirical expectations based on the
information content of behavioral experiments and the degrees of freedom in
the $\Psi$-space. The estimates assume proper regularization and
cross-validation; overfitting would yield high in-sample but low
out-of-sample $R^2$.

%==============================================================================
% PART V: ANTICIPATED OBJECTIONS
%==============================================================================

\subsubsection{Anticipated Objections}

A skeptical reader---particularly one trained in experimental methods---will
have immediate concerns. We address them directly.

\paragraph{``Isn't this just curve-fitting?''}
No. Curve-fitting would mean: train on experiment $X$, predict experiment $X$.
The test is: train on experiments $X_1, \ldots, X_n$ (e.g., Western ultimatum
games), predict experiment $Y$ (e.g., Machiguenga ultimatum games) that was
held out from training.

If the model captures the $\Psi$-dependence correctly, it generalizes. If it
has merely memorized data points, it fails on held-out experiments. This is
the standard logic of cross-validation, applied to behavioral theory.

\paragraph{``How do you avoid overfitting with so many parameters?''}
Three mechanisms:
\begin{enumerate}
    \item \textbf{Regularization}: Ridge/LASSO penalties shrink coefficients
          toward zero unless supported by data
    \item \textbf{Cross-validation}: Parameters are chosen to minimize
          out-of-sample error, not in-sample fit
    \item \textbf{Structural constraints}: The 8 $\Psi$-dimensions are not
          arbitrary---they are the dimensions along which experimental
          evidence shows behavior varies. Adding a 9th dimension requires
          passing the Type III diagnostic test (Appendix~\ref{app:epistemics}):
          the candidate must explain residual variance not captured by existing
          dimensions.
\end{enumerate}

\paragraph{``Why should I believe your $\Psi$ measurements?''}
You shouldn't---uncritically. The $\Psi$ values come from established
instruments: World Values Survey for $\Psi_S$ and cultural values,
Global Preference Survey for time/risk preferences, V-Dem for institutional
quality. Each has known measurement properties, validated across contexts.

But any measurement error in $\Psi$ shows up as residual variance in
$C^*(\Psi)$. The diagnostic protocol (Appendix~\ref{app:epistemics})
classifies this: Type I deviations indicate measurement noise; Type III
deviations suggest missing dimensions. The framework does not assume perfect
measurement---it diagnoses imperfection.

\paragraph{``What about experimenter demand effects, publication bias, and other artifacts?''}
These are real concerns for any meta-analytic approach. We address them through:
\begin{itemize}
    \item \textbf{Design coding}: Each experiment is coded for double-blind
          status, incentive compatibility, and other design features that
          moderate effect sizes
    \item \textbf{Heterogeneity modeling}: Publication bias inflates average
          effect sizes but not their $\Psi$-dependence---if anything, it
          attenuates moderator effects toward zero
    \item \textbf{Robustness checks}: Key predictions are tested against
          subsets excluding controversial or low-quality studies
\end{itemize}

\paragraph{``Can this really work for all of behavioral economics?''}
No. The framework has scope conditions:
\begin{itemize}
    \item \textbf{In scope}: Behaviors where complementarity and context
          matter---social preferences, risk, time, cooperation, fairness,
          trust, reciprocity
    \item \textbf{Partially in scope}: Cognitive biases that interact with
          context (framing effects, defaults)
    \item \textbf{Out of scope}: Pure perceptual or memory errors that do not
          depend on social/institutional context
\end{itemize}

The claim is not ``EBF explains everything.'' It is: ``For behaviors
that depend on complementarity and context, EBF provides a calibrated
predictive model that LLMs cannot match.''

%==============================================================================
% PART VI: THE EXPERIMENTAL RECORD AS SCIENTIFIC INFRASTRUCTURE
%==============================================================================

\subsubsection{The Experimental Record as Scientific Infrastructure}

The behavioral economics literature represents an extraordinary scientific
resource: thousands of experiments, conducted over 50 years, across dozens
of countries, probing every aspect of human decision-making.

But this resource has been underutilized. Each paper stands alone. Meta-analyses
aggregate within narrow domains. No framework has systematically extracted
the \emph{functional relationships} encoded in this evidence.

\paragraph{What Full Calibration Requires.}
To realize the potential of this literature, each experiment must be coded
not just for its headline result, but for:

\begin{itemize}
    \item \textbf{Context ($\Psi$)}: Country, population, institutional setting,
          with values from independent measurement instruments
    \item \textbf{Design ($X$)}: Stakes (real/hypothetical, amount), framing,
          repetition, matching protocol, anonymity, sample characteristics
    \item \textbf{Results ($Y$)}: Full distributions where available, not just
          means; effect sizes with standard errors; treatment-control contrasts
    \item \textbf{Moderators ($M$)}: What variations were tested? What
          interactions were found or not found?
\end{itemize}

With this information, the experimental record becomes a \emph{training set}
for $C^*(\Psi)$---not in the sense of machine learning on text, but in the
sense of scientific calibration on behavioral data.

\paragraph{The BEATRIX Implementation.}
The BEATRIX system operationalizes this vision:
\begin{itemize}
    \item A structured database of experimental findings with full metadata
    \item Automated extraction of $\Psi$-values from context descriptions,
          linked to WVS/GPS/V-Dem measurements
    \item Continuous updating of $C^*(\Psi)$ as new experiments are coded
    \item Diagnostic protocols that flag when predictions fall outside
          the calibration range or when residuals suggest model revision
    \item Version control: every change to $C^*(\Psi)$ is logged with the
          evidence that motivated it
\end{itemize}

This is not ``artificial intelligence'' in the sense of LLMs. It is
\emph{scientific infrastructure} that makes the accumulated evidence of
behavioral economics computationally accessible.

%==============================================================================
% PART VII: THE TESTABLE CLAIM
%==============================================================================

\subsubsection{The Testable Claim}

The argument of this section reduces to a falsifiable prediction:

\begin{center}
\fbox{\parbox{0.9\textwidth}{
\textbf{Prediction}: A EBF model calibrated on $N$ experiments can predict
the outcomes of held-out experiments (probing the same mechanisms but not
included in calibration) with accuracy exceeding:
\begin{enumerate}
    \item[(a)] Na\"ive baselines (population means, simple demographics)
    \item[(b)] LLM-based predictions (GPT-4, Claude, Gemini)
    \item[(c)] Expert predictions (behavioral economists forecasting outcomes)
\end{enumerate}
The accuracy advantage increases with $N$ and with the diversity of the
calibration set.
}}
\end{center}

This prediction can be tested today:
\begin{enumerate}
    \item Assemble a calibration set of $N$ experiments with full coding
    \item Hold out 10--20\% as a test set, stratified by domain and context
    \item Calibrate $C^*(\Psi)$ on the training set
    \item Generate predictions for the test set (with confidence intervals)
    \item Elicit predictions from (a) baselines, (b) LLMs, (c) experts
    \item Compare using proper scoring rules (Brier score, log loss)
\end{enumerate}

If EBF outperforms all comparators, the framework is validated---not as
a theoretical construct, but as a \emph{predictive engine} for human behavior.

If it does not, the framework is refuted, or at minimum, its scope conditions
are identified.

Either outcome advances science. But the prediction of this paper is that
EBF, properly calibrated, will outperform alternatives---because it is
built on the right foundation: not text about behavior, but behavioral data
themselves.

%==============================================================================
% PART VIII: IMPLICATIONS FOR THE DIGITAL TWIN DEBATE
%==============================================================================

\subsubsection{Implications for the Digital Twin Debate}

The \citet{peng2025} findings generated considerable discussion about the
future of LLM-based behavioral simulation. Our analysis suggests a resolution:

\paragraph{Why Digital Twins Fail.}
Digital twins fail not because LLMs are insufficiently powerful, but because
they lack a \emph{model} of behavior. They interpolate from text, which
encodes descriptions of behavior, not the behavioral regularities themselves.

The problem is architectural. No amount of scaling---larger models, more
training data, better prompts---will solve it, because the training data
(text) does not contain the information needed for prediction (quantified
functional relationships between context and behavior).

\paragraph{What Would Make Them Work.}
For LLM-based simulation to succeed at behavioral prediction, it would need:
\begin{enumerate}
    \item Explicit representation of context ($\Psi$)
    \item Functional relationships between context and behavior ($C^*(\Psi)$)
    \item Parameters estimated from behavioral data, not inferred from text
    \item Uncertainty quantification for predictions outside training range
\end{enumerate}

But this is precisely what EBF provides---without requiring LLMs at all.
The ``intelligence'' is not in the language model; it is in the
\emph{calibrated behavioral model}.

\paragraph{The Role of LLMs in BEATRIX.}
This does not mean LLMs are useless. In the BEATRIX architecture, LLMs serve
specific functions:
\begin{itemize}
    \item \textbf{Natural language interface}: Translating user queries into
          $\Psi$-space specifications
    \item \textbf{Report generation}: Presenting predictions and diagnostics
          in accessible form
    \item \textbf{Literature processing}: Extracting experimental details from
          papers for database coding
    \item \textbf{Hypothesis generation}: Suggesting experimental designs to
          test model boundaries
\end{itemize}

But LLMs do not generate the behavioral predictions. $C^*(\Psi)$ does---
calibrated on the accumulated evidence of experimental behavioral economics.

%==============================================================================
% PART IX: FORWARD REFERENCE
%==============================================================================

\paragraph{Forward Reference.}
The predictions generated by this calibration approach are developed in
Chapter~\ref{sec:novelpredictions}. There we show that $C^*(\Psi)$ generates
predictions about trade wars, technological disruption, institutional reform,
and other policy-relevant phenomena that:
\begin{itemize}
    \item No single constituent theory produces alone
    \item Differ systematically from both standard economic predictions and
          LLM-based forecasts
    \item Can be tested empirically with pre-registered observational and
          experimental designs
\end{itemize}

The calibration logic developed in this section is what makes those predictions
possible---and credible.

%==============================================================================
% CONCLUSION
%==============================================================================

\vspace{1em}
\begin{center}
\fbox{\parbox{0.92\textwidth}{
\textbf{The Bottom Line}: Fifty years of behavioral economics experiments
are not just evidence that humans deviate from rationality. They are the
\emph{calibration data} for a predictive model of those deviations. No
language model, however large, can replicate this---because prediction
requires data about behavior, not text about behavior.

\vspace{0.5em}

The experiments of Fehr, Kahneman, Thaler, and hundreds of others did not
just show that people exhibit social preferences, loss aversion, or
present bias. They showed \emph{how these phenomena vary with context}.
That variation is the signal. And EBF is the first framework to
systematically extract it.

\vspace{0.5em}

This is why BEATRIX can replicate experiments that digital twins cannot.
This is why its accuracy scales with the experimental record. And this is
why---we predict---it will outperform any LLM on behavioral forecasting.

\vspace{0.5em}

The test is straightforward. The prediction is falsifiable. And the stakes
are clear: either behavioral economics has produced cumulative knowledge
that can be computationally leveraged, or it has not.

\vspace{0.5em}

We believe it has.
}}
\end{center}

%%%%%%%%%%%%%%%%%%%%%%%%%%%%%%%%%%%%%%%%%%%%%%%%%%%%%%%%%%%%%%%%%%%%%%%%%%%%%%%
% BIBLIOGRAPHY ENTRIES (to be added to main bibliography)
%%%%%%%%%%%%%%%%%%%%%%%%%%%%%%%%%%%%%%%%%%%%%%%%%%%%%%%%%%%%%%%%%%%%%%%%%%%%%%%
%
% @article{peng2025,
%   title={A Mega-Study of Digital Twins Reveals Strengths, Weaknesses and 
%          Opportunities for Further Improvement},
%   author={Peng, Tianyi and Gui, George and Merlau, Daniel J. and others},
%   journal={arXiv preprint arXiv:2509.19088},
%   year={2025}
% }
%
% @techreport{horton2023,
%   title={Large Language Models as Simulated Economic Agents: What Can We Learn 
%          from Homo Silicus?},
%   author={Horton, John J.},
%   year={2023},
%   institution={National Bureau of Economic Research},
%   number={31122}
% }
%
% @article{argyle2023,
%   title={Out of One, Many: Using Language Models to Simulate Human Samples},
%   author={Argyle, Lisa P. and Busby, Ethan C. and Fulda, Nancy and Gubler, 
%           Joshua R. and Rytting, Christopher and Wingate, David},
%   journal={Political Analysis},
%   volume={31},
%   number={3},
%   pages={337--351},
%   year={2023}
% }
%
% @inproceedings{park2023,
%   title={Generative Agents: Interactive Simulacra of Human Behavior},
%   author={Park, Joon Sung and O'Brien, Joseph C. and Cai, Carrie J. and 
%           Morris, Meredith Ringel and Liang, Percy and Bernstein, Michael S.},
%   booktitle={Proceedings of the 36th Annual ACM Symposium on User Interface 
%              Software and Technology},
%   year={2023}
% }
%
% @article{henrich2001,
%   title={In Search of Homo Economicus: Behavioral Experiments in 15 
%          Small-Scale Societies},
%   author={Henrich, Joseph and Boyd, Robert and Bowles, Samuel and 
%           Camerer, Colin and Fehr, Ernst and Gintis, Herbert and others},
%   journal={American Economic Review},
%   volume={91},
%   number={2},
%   pages={73--78},
%   year={2001}
% }
%
% @article{guth1982,
%   title={An Experimental Analysis of Ultimatum Bargaining},
%   author={G{\"u}th, Werner and Schmittberger, Rolf and Schwarze, Bernd},
%   journal={Journal of Economic Behavior \& Organization},
%   volume={3},
%   number={4},
%   pages={367--388},
%   year={1982}
% }
%
% @article{johnson2003,
%   title={Do Defaults Save Lives?},
%   author={Johnson, Eric J. and Goldstein, Daniel},
%   journal={Science},
%   volume={302},
%   number={5649},
%   pages={1338--1339},
%   year={2003}
% }
%


%%%%%%%%%%%%%%%%%%%%%%%%%%%%%%%%%%%%%%%%%%%%%%%%%%%%%%%%%%%%%%%%%%%%%%%%%%%%%%%
% SUMMARY SECTION
%%%%%%%%%%%%%%%%%%%%%%%%%%%%%%%%%%%%%%%%%%%%%%%%%%%%%%%%%%%%%%%%%%%%%%%%%%%%%%%
\subsection{Summary}
\label{sec:ch4x-summary}

This chapter established the fundamental distinction between \emph{calibration} and \emph{simulation}:

\begin{enumerate}
\item \textbf{LLM simulation fails}: Digital twins achieve only 75\% individual-level accuracy---no better than demographic baselines. They under-represent human heterogeneity and lack the phenomenological grounding needed for behavioral prediction.

\item \textbf{Calibration succeeds}: The EBF approach uses LLMs to extract parameters from 50+ years of experimental research, then applies these parameters in theoretically grounded models with clear epistemic status tags.

\item \textbf{Experimental replication is guaranteed}: Because parameters come from replicated experiments, predictions inherit their reliability. The 1,500+ studies in the EBF database provide robust calibration across domains.

\item \textbf{The key insight}: LLMs are pattern recognizers, not experience simulators. Using them for calibration leverages their strength; using them for simulation exposes their weakness.
\end{enumerate}

\textbf{Formal Details:} Complete estimation methodology, epistemic status tags (EMP/THR/LLM/ILL/HYP), and validation protocols are provided in Appendix BBB (CORE-WHERE).

\textbf{What Comes Next:} Chapter 5 introduces the complementarity structure $\gamma$ that these calibrated parameters populate.

%%%%%%%%%%%%%%%%%%%%%%%%%%%%%%%%%%%%%%%%%%%%%%%%%%%%%%%%%%%%%%%%%%%%%%%%%%%%%%%
% READING PATH BOX
%%%%%%%%%%%%%%%%%%%%%%%%%%%%%%%%%%%%%%%%%%%%%%%%%%%%%%%%%%%%%%%%%%%%%%%%%%%%%%%
\begin{tcolorbox}[
    colback=green!5!white,
    colframe=green!75!black,
    title={\textbf{Reading Path}},
    fonttitle=\bfseries
]
\begin{tabular}{@{}ll@{}}
\textbf{Previous:} & Chapter 4 (Empirical Foundations) \\
\textbf{Next:} & Chapter 5 (Complementarity) \\
\textbf{Deep Dive:} & See Appendix BBB (Parameter Repository) for complete estimation methodology \\
\end{tabular}

\medskip
\textit{For readers short on time: Focus on the Intuition Box and Summary section.}
\end{tcolorbox}



%------------------------------------------------------------------------------
% SUMMARY
%------------------------------------------------------------------------------
\section{Summary}
\label{sec:ch4-summary}

\begin{tcolorbox}[
    colback=blue!5!white,
    colframe=blue!75!black,
    title={\textbf{Chapter 4 Summary: Empirical Foundations}}
]

\textbf{Core Insight:} The EBF framework requires calibration on experimental data, not derivation from axioms or LLM-based simulation.

\textbf{Key Results:}
\begin{enumerate}[nosep]
    \item \textbf{Parameter Estimation:} Key behavioral parameters ($\lambda$, $\beta$, $\alpha$) come from 1,900+ experimental papers
    \item \textbf{Epistemic Tags:} EMP (empirical), THR (theoretical), LLM (estimated) classify evidence quality
    \item \textbf{Calibration vs Simulation:} LLMs reproduce words but not phenomena---calibration required
    \item \textbf{Continuous Learning:} Self-reinforcement architecture updates parameters from new evidence
\end{enumerate}

\textbf{Transition:} Chapter 5 develops complementarity ($\gamma$) as the structural principle governing utility interactions.
\end{tcolorbox}

\subsection{What Comes Next}
\label{sec:ch4-next}

With the empirical foundations established, subsequent chapters build the theoretical structure:

\begin{enumerate}[nosep]
    \item \textbf{Chapter 5:} Introduces complementarity as the structural principle (CORE-HOW)
    \item \textbf{Chapter 6:} Derives the reference structure $C^*(\Psi)$
    \item \textbf{Chapter 9:} Develops context as an endogenous variable (CORE-WHEN)
\end{enumerate}

%%%%%%%%%%%%%%%%%%%%%%%%%%%%%%%%%%%%%%%%%%%%%%%%%%%%%%%%%%%%%%%%%%%%%%%%%%%%%%%
% READING PATH BOX
%%%%%%%%%%%%%%%%%%%%%%%%%%%%%%%%%%%%%%%%%%%%%%%%%%%%%%%%%%%%%%%%%%%%%%%%%%%%%%%
\begin{tcolorbox}[
    colback=green!5!white,
    colframe=green!75!black,
    title={\textbf{Reading Path}},
    fonttitle=\bfseries
]
\begin{tabular}{@{}ll@{}}
\textbf{Previous:} & Chapter 3 \\
\textbf{Next:} & Chapter 5 \\
\textbf{Deep Dive:} & See Appendix G (Glossary) for definitions \\
\textbf{Parameterization:} & Appendix UNMAPPED_PM (complete methodology navigation) \\
\end{tabular}

\medskip
\textit{For readers short on time: Focus on the Intuition Box and Summary section.}
\end{tcolorbox}

